% concatenated from arxiv.tex
\documentclass{article}

% if you need to pass options to natbib, use, e.g.:
%     \PassOptionsToPackage{numbers, compress}{natbib}
% before loading neurips_2026

% The authors should use one of these tracks.
% Before accepting by the NeurIPS conference, select one of the options below.
% 0. "default" for submission
\usepackage[preprint, nonatbib]{neurips_2026}
\usepackage{longtable}
\usepackage{booktabs}
% the "default" option is equal to the "main" option, which is used for the Main Track with double-blind reviewing.
% 1. "main" option is used for the Main Track
%  \usepackage[main]{neurips_2026}
% 2. "position" option is used for the Position Paper Track
%  \usepackage[position]{neurips_2026}
% 3. "eandd" option is used for the Evaluations & Datasets Track
 % \usepackage[eandd]{neurips_2026}
 % if you need to opt-in for a single-blind submission in the E&D track:
 %\usepackage[eandd, nonanonymous]{neurips_2026}
% 4. "creativeai" option is used for the Creative AI Track
%  \usepackage[creativeai]{neurips_2026}
% 5. "sglblindworkshop" option is used for the Workshop with single-blind reviewing
 % \usepackage[sglblindworkshop]{neurips_2026}
% 6. "dblblindworkshop" option is used for the Workshop with double-blind reviewing
%  \usepackage[dblblindworkshop]{neurips_2026}

% After being accepted, the authors should add "final" behind the track to compile a camera-ready version.
% 1. Main Track
 % \usepackage[main, final]{neurips_2026}
% 2. Position Paper Track
%  \usepackage[position, final]{neurips_2026}
% 3. Evaluations & Datasets Track
 % \usepackage[eandd, final]{neurips_2026}
% 4. Creative AI Track
%  \usepackage[creativeai, final]{neurips_2026}
% 5. Workshop with single-blind reviewing
%  \usepackage[sglblindworkshop, final]{neurips_2026}
% 6. Workshop with double-blind reviewing
%  \usepackage[dblblindworkshop, final]{neurips_2026}
% Note. For the workshop paper template, both \title{} and \workshoptitle{} are required, with the former indicating the paper title shown in the title and the latter indicating the workshop title displayed in the footnote.
% For workshops (5., 6.), the authors should add the name of the workshop, "\workshoptitle" command is used to set the workshop title.
% \workshoptitle{WORKSHOP TITLE}

% "preprint" option is used for arXiv or other preprint submissions
 % \usepackage[preprint]{neurips_2026}

% to avoid loading the natbib package, add option nonatbib:
%    \usepackage[nonatbib]{neurips_2026}

\usepackage[utf8]{inputenc} % allow utf-8 input
\usepackage[T1]{fontenc}    % use 8-bit T1 fonts
\usepackage{hyperref}       % hyperlinks
\usepackage{url}            % simple URL typesetting
\usepackage{booktabs}       % professional-quality tables
\usepackage{amsfonts}       % blackboard math symbols
\usepackage{nicefrac}       % compact symbols for 1/2, etc.
\usepackage{microtype}      % microtypography
\usepackage{xcolor}         % colors
\usepackage{cite}
\usepackage{multirow}
\usepackage{booktabs}
\usepackage{wrapfig}

% Note. For the workshop paper template, both \title{} and \workshoptitle{} are required, with the former indicating the paper title shown in the title and the latter indicating the workshop title displayed in the footnote. 
% \title{Inference-Time Alignment of Video Generation Models via Control-Theoretic Activation Steering}
\title{Activation Steering of Video Generation Models via Reduced-Order Linear Optimal Control}


\input{preamble}



\author{%
  Jihoon Hong\qquad Alice Chan$^\star$\qquad Qiyue Dai$^\star$\qquad Julian Skifstad\qquad Glen Chou\\ %\thanks{Use footnote for providing further information about author (webpage, alternative address)---\emph{not} for acknowledging funding agencies.} \\
  Georgia Institute of Technology\\
  Atlanta, GA 30308 \\
  \texttt{\{jhong392, ichan30, qdai41, jskifstad3, chou\}@gatech.edu} \\
  % examples of more authors
  % \And
  % Coauthor \\
  % Affiliation \\
  % Address \\
  % \texttt{email} \\
  % \AND
  % Coauthor \\
  % Affiliation \\
  % Address \\
  % \texttt{email} \\
  % \And
  % Coauthor \\
  % Affiliation \\
  % Address \\
  % \texttt{email} \\
  % \And
  % Coauthor \\
  % Affiliation \\
  % Address \\
  % \texttt{email} \\
}


\begin{document}


\maketitle

\begin{abstract}
\looseness-1Text-to-video (T2V) models trained on large-scale web data can generate undesired content, motivating interventions that reduce harmful outputs without sacrificing visual quality. Activation steering offers an attractive mechanistic alternative to finetuning and prompt filtering, but existing T2V steering methods remain limited, typically applying coarse, non-anticipative interventions that can lead to oversteering and content degradation. 
To close this gap, we propose Latent Activation Linear-Quadratic Regulator (LA-LQR), a reduced-order optimal control framework for minimally invasive T2V steering. LA-LQR formulates T2V inference as a dynamical system and computes closed-loop feedback interventions that steer activations toward desired feature setpoints while penalizing unnecessary perturbations. To make optimal control feasible for high-dimensional video activations, we project activations onto a low-dimensional, task-relevant subspace derived from contrastive prompt pairs, estimate local linear dynamics in this latent space, and solve a latent LQR problem to obtain timestep- and layer-specific steering signals. We provide theoretical bounds relating latent setpoint tracking to raw activation-space feature control, and empirically validate the fidelity of the reduced latent dynamics. On concept steering and video safety benchmarks, LA-LQR reduces unsafe generations relative to baselines, while preserving prompt fidelity and visual quality.
\end{abstract}

\begin{center}
\vspace{-5pt}
    \textcolor{red}{Warning: This paper contains offensive model outputs.}
\end{center}

\vspace{-15pt}
\section{Introduction}
\vspace{-5pt}

Text-to-video (T2V) models \cite{wan2025wan, ju2025editverse, zheng2024open, yang2024cogvideox, ho2022video} can generate high-fidelity videos from text, with many applications \cite{wu2025moviebench, he2026pre, kim2026cosmos}. However, because they are trained on weakly curated web-scale data, they may inherit undesired concepts, including nudity, graphic violence, copyrighted content, and depictions of public figures \cite{facchiano2025video, yoon2024safree}. This raises serious safety risks, including deepfake-enabled misinformation, abuse, and fraud, motivating methods that reliably constrain outputs while preserving model utility.

\looseness-1Recent work aims to improve T2V model safety. Model finetuning can suppress harmful concepts, but requires extensive training resources \cite{cheng2025vpo, dai2024safesora}. Model editing methods~\cite{gandikota2024unified} directly modify weights to improve safety, but can degrade output quality and make models behave unpredictably. Filtering methods reject or rewrite unsafe prompts without altering the model, but are vulnerable to jailbreaks and indirect prompts \cite{yoon2024safree}. Inference-time alignment methods, such as activation steering~\cite{turner2024activation, ode_activation_steering, skifstad2026local}, offer a promising alternative because they can mechanistically modulate model behavior without retraining or altering weights. However, existing methods are often \textit{non-anticipative}: they apply fixed or myopic interventions without accounting for how perturbations propagate through the model. As a result, they can over- or under-steer generations, reduce prompt fidelity, or remain sensitive to adversarial prompts.

\looseness-1To address these limitations, we propose a mechanistic reduced-order optimal control framework, called \underline{L}atent \underline{A}ctivation \underline{L}inear-\underline{Q}uadratic \underline{R}egulator (LA-LQR), for minimally invasive T2V steering. We formalize T2V generation as a dynamical system governed by the model’s weights and architecture, with steering signals as control inputs. This yields a controller that reaches desired concept-modulation setpoints while minimizing perturbations, reducing oversteering and preserving video quality. Using local linearizations of T2V dynamics, LA-LQR anticipates natural activation evolution and avoids unnecessary interventions when activations already move toward the target. Online error feedback further improves robustness by adjusting steering based on deviations from the setpoint. Since full activation-space optimal control is infeasible for T2V models with tens of millions of activation dimensions, we project activations onto a low-dimensional, task-relevant latent feature subspace derived from contrastive prompt pairs. In this space, we estimate local linear dynamics with efficient Jacobian-vector products, then compute minimum-norm interventions toward desired latent feature setpoints. We validate the accuracy and concept fidelity of the resulting linear latent dynamics and demonstrate a reduction in unsafe outputs while preserving prompt fidelity and visual quality. Our contributions are:
\vspace{-7pt}

\begin{wrapfigure}{r}{.58\textwidth}
\centering
\vspace{-28pt}
    \includegraphics[width=\linewidth]{figs/front_fig_v2.pdf}\vspace{-21pt}
    \caption{\textbf{Overview}. Our method, LA-LQR, steers T2V models by solving an optimal control problem, producing steering signals $u_s$ for each timestep and transformer layer. For tractability, we perform control within a task-relevant activation subspace identified by contrastive vectors.\vspace{-10pt}}
    \vspace{-3pt}
    \label{fig:overview}
\end{wrapfigure}

\begin{itemize}
    \item We \textit{formalize T2V generation} as a finite-horizon dynamical system 
    and propose a \textit{linear optimal control framework for inference-time T2V steering}, deriving minimally invasive online feedback interventions on text- and video-token activations.
    
    \item We make control tractable for T2V models by constructing and projecting to a \textit{reduced-order latent space} capturing dominant features from contrastive vectors in the full activation space.
    
    \item We \textit{justify the reduced-order controller theoretically and empirically} by quantifying the feature information lost in the projection and validating local linearity in the latent space.
    
    \item We \textit{evaluate on concept steering and video-safety tasks}, reducing unsafe content while preserving prompt fidelity and visual quality, surpassing T2V steering and safety baselines.
\end{itemize}


\vspace{-6pt}
\section{Related Work}
\vspace{-4pt}

\paragraph{Activation Steering}
\looseness-1Prior work in mechanistic interpretability \cite{bereska2024mechanistic, Elhage_Hume_Olsson_2022, mikolov2013linguistic, park2023linear, marks2024the, nanda2023emergent} suggests that many concepts align with directions in activation space, enabling behavior modulation along those directions. Activation steering builds on this insight by perturbing activations during inference, primarily in LLMs. Most methods use contrastive addition \cite{Dathathri2020Plug}, deriving steering vectors from examples with and without a target concept \cite{li2023inference, turner2024activation, Arditi_Obeso_Syed_Paleka_Panickssery_Gurnee_Nanda_2024a, Rimsky_Gabrieli_Schulz_Tong_Hubinger_Turner_2024, Li_Patel_Viégas_Pfister_Wattenberg_2024, Turner_Thiergart_Leech_Udell_Vazquez_Mini_MacDiarmid_2024}. Although this can offer more systematic alignment than finetuning \cite{stiennon2020learning, touvron2023llama, rafailov2023direct, xu2024contrastive, yuan2023rrhf, song2024preference, houlsby2019parameter}, prompting \cite{askell2021general, zhang2023defending}, or guided decoding \cite{khanov2024args, huang2025deal}, many methods use simple interventions \cite{Turner_Thiergart_Leech_Udell_Vazquez_Mini_MacDiarmid_2024, Rodriguez_Blaas_Klein_Zappella_Apostoloff_Cuturi_Suau_2024, wu2024reft, wu2024advancing, Vu_Nguyen_2025} without modeling perturbation propagation, which can lead to oversteering. 
Control-theoretic steering methods address this by treating LLMs as dynamical systems \cite{bhargava2023s}, but often rely on offline controller training \cite{Kong_Wang_Mu_Du_Zhuang_Zhou_Song_Zhang_Wang_Zhang_2024, Karnik_Bansal_2025, Cheng_Alonso_2025, Hedstrom_Amoukou_Bewley, tan2024analysing, scalena2024multi, ode_activation_steering, miyaoka2024cbf, Nguyen_Vu_Pham_Zhang_Nguyen_2025}. \cite{skifstad2026local} uses local linearity to derive an optimal controller for minimally invasive online LLM steering. However, it does not apply to diffusion transformer models with spatiotemporal denoising, and T2V models’ much higher activation dimensionality makes the full-space control problem of \cite{skifstad2026local} intractable, requiring petabytes of memory. We extend this control-theoretic framing to T2V models and introduce a \textit{reduced-order} linear control formulation for feasible online feedback steering. 


\vspace{-10pt}
\paragraph{T2V Attacks \& Defenses}
T2V models \cite{cerspense2023zeroscope, hong2023cogvideo, yang2025cogvideox, wan2025wan, ho2022imagenvideohighdefinition, singer2022makeavideotexttovideogenerationtextvideo, wang2023modelscope, liu2024sorareviewbackgroundtechnology, opensora} can learn unsafe concepts that adversaries can exploit. Prior work constructs prompt-level attacks~\cite{li2025t2vattack, liu2025t2v, miao2024t2vsafetybench} that can evade frame-based moderation~\cite{chen2026two, Wang_2025_ICCV}. Existing defenses use filtering, prompt-time defenses, or concept erasure: \cite{yoon2024safree} masks unsafe prompt tokens, \cite{liu2024unlearning} uses trigger-token embeddings, \cite{yi5993786nullsce} combines negative noise guidance with diffuser fine-tuning, and \cite{facchiano2025video} applies low-rank weight updates. These methods either act mainly at the input-token level, are vulnerable to jailbreaks, require retraining, or lack inference-time flexibility.
For activation steering, most defenses target LLMs, with limited extensions to image generation~\cite{Rodriguez_Blaas_Klein_Zappella_Apostoloff_Cuturi_Suau_2024, rodriguez2025lineas, Nguyen_Vu_Pham_Zhang_Nguyen_2025} but not T2V. For T2V, \cite{facchiano2025video} applies weight modifications that induce activation updates resembling \cite{turner2024activation}. A recent work \cite{ekin2026unreasonable} performs activation steering for T2V models.
However, since T2V models use text-encoder outputs as context~\cite{raffel2020exploring}, this distributes semantics across text token embeddings and makes LLM-style last-token steering \cite{Karnik_Bansal_2025} inapplicable. Consequently, \cite{ekin2026unreasonable} requires an LLM to select embedding positions, which is unreliable for jailbroken safety prompts. In contrast, we show in Sec.~\ref{sec:results} that online, minimally invasive closed-loop feedback using optimal control enables reliable T2V content modulation without LLM-assisted token selection.

\vspace{-7pt}
\section{Preliminaries and Problem Statement}
\vspace{-4pt}

\paragraph{T2V Architectures}
We consider T2V models based on Diffusion Transformers (DiTs) \cite{Peebles2022DiT}, which are trained to learn a $T$-step iterative reverse diffusion process that can be generally described as:
\begin{equation}
    \hat x_T \sim \mathcal{N}(0,I),\qquad
    h = \texttt{Embed}(p),
\end{equation}
\begin{equation}
    \begin{aligned}
        \mu_{t} = M(\hat x_t, h, t), \qquad
        \hat x_{t-1} = \texttt{Scheduler}(\hat x_t, \mu_t, t), \qquad
        x_{\text{out}} = \texttt{Dec}(\hat x_1),
    \end{aligned}
\end{equation}

where $\mathcal{N}(0,I)$ is a unit Gaussian, $\hat x_t$ is the latent video representation, $p$ is the text prompt, $\texttt{Embed}$ is the text embedder, $\texttt{Scheduler}$ is a reverse diffusion (or flow) scheduler, $\texttt{Dec}$ is the decoder, $x_\textrm{out}$ is the output video, and $M$ is an $L$-layer DiT, where we omit explicit timestep conditioning for brevity:
\begin{subequations}\label{eq:transformer}
\begin{align}
    M &:= \texttt{Unpatchify} \circ \phi^{(L)} \circ ... \circ \phi^{(1)} \circ \texttt{Patchify}(\hat x_t)\label{eq:transformer_M} \\
    x_{l+1,t} &:= \phi^{(l)}(x_{l,t}) := x_{\text{attn}} + \texttt{MLP}_l(\texttt{Norm}(x_{\text{attn}})), \quad x_{1,t}:= \texttt{Patchify}(\hat x_t),\label{eq:transformer_dynamics} \\
    x_{\text{attn}} &:= x' + \texttt{CrossAttn}_l(\texttt{Norm}(x'), h), \qquad
    x' := x_{l,t} + \texttt{SelfAttn}_l(\texttt{Norm}(x_{l,t}))\label{eq:transformer_z},
\end{align}
\end{subequations}

For a video with $F$ frames, $Q$ patches per frame, and hidden dimension $d$, the latent video representation at layer $l$ and timestep $t$ is $x_{l,t} \in \mathbb{R}^{F \times Q \times d}$. For a prompt of length $P$ and text hidden dimension $d'$, the text representation is $h \in \mathbb{R}^{P \times d'}$. A maximum context length $C$ is often fixed, so $h \in \mathbb{R}^{C \times d'}$.

\vspace{-8pt}
\paragraph{Linear Quadratic Regulator (LQR)}\label{sec:prelim_lqr}

We adapt LQR \cite{kalman1960contributions} to compute steering policies for T2V models. LQR considers an optimal control problem for linear time varying (LTV) dynamics \eqref{eq:ltv}:
\begin{subequations}\label{eq:lqr}
\begin{align}
\textstyle\min_{\{u_k\}_{k=1}^{H-1}}\quad
& J := z_H^\top Q_H z_H
+ \textstyle \sum_{k=1}^{H-1} \left( z_k^\top Q_k z_k + u_k^\top R_k u_k \right)
\label{eq:lqr_objective} \\
\text{subject to} \quad\ \
& z_{k+1} = A_k z_k + B_k u_k, \qquad \forall k \in \{1,\dots,H\} := [H-1],
\label{eq:ltv}
\end{align}
\end{subequations}
\looseness-1for state $z_k$ and control $u_k$, where $A_k,B_k,Q_k,R_k$ are appropriately-sized, $Q_k \succeq 0,\ \forall k \in [H]$ and $R_k \succ 0,\ \forall k \in [H-1]$. LQR returns an optimal $\{u^*_k\}_{k=1}^{H-1}$ that minimizes $J$, yielding a closed-form solution $u^*_k = -K_kz_k$, where the gain $K_k$ is found via Riccati recursions 
\cite{lewisOptimalControlChapter2}. The objective \eqref{eq:lqr_objective} penalizes deviations from $(\bar z,\bar u) = (0,0)$ and can be generalized to penalize deviation from desired setpoints $(\{\bar z_k\}_{k=1}^H, \{\bar u_k\}_{k=1}^{H-1})$. Denoting $\delta z_k := z_k - \bar z_k$ and $\delta u_k := u_k - \bar u_k$, we can modify \eqref{eq:lqr} as
\begin{subequations}\label{eq:lqr_tracking}
\begin{align}
\textstyle\min_{\{\delta u_k\}_{k=1}^{H-1}}\quad
& \delta z_H^\top Q_H \delta z_H
+ \textstyle\sum_{k=1}^{H-1} \left( \delta z_k^\top Q_k \delta z_k + \delta u_k^\top R_k \delta u_k \hspace{-1pt}\right)
\label{eq:err_lqr_objective} \\
\text{subject to} \quad\ \
& \delta z_{k+1} = A_k \delta z_k + B_k \delta u_k, \qquad k = 1,\dots,H-1,
\label{eq:err_ltv}
\end{align}
\end{subequations}
to compute a corresponding optimal trajectory tracking controller 
% \begin{equation}\label{eq:lqr_u}
$u^*_k := \bar u_k +\delta u^*_k := \bar u_k -K_k\delta z_k$.
% \end{equation}
LQR can be applied to nonlinear systems $z_{k+1} = f_k(z_k, u_k)$ via first order Taylor expansions about a nominal trajectory $\{(\bar z_k, \bar u_k)\}_{k=1}^T$, yielding the approximation
% we can approximate 
\begin{equation}\label{lqr_linearization}
\begin{aligned}
    \bar z_{k+1} + \delta z_{k+1} &= f_k(\bar z_k + \delta z_k, \bar u_k + \delta u_k) \approx f_k(\bar z_k,\bar u_k) + A_k\delta z_k + B_k\delta u_k, \\ 
    \delta z_{k+1} &\approx A_k\delta z_k + B_k\delta u_k,
\end{aligned}
\end{equation}
for Jacobians $A_k:=(\partial f_k/\partial z)|_{z_k,u_k}$, $B_k := (\partial f_k/\partial u)|_{z_k,u_k}$. We can then define an LQR problem analogous to \eqref{eq:lqr_tracking} to find a controller penalizing deviations from the nominal trajectory $\{(\bar z_k, \bar u_k)\}_{k=1}^T$.


\vspace{-8pt}
\paragraph{Problem Statement}
\looseness-1We control a T2V model by intervening on internal activations to modulate an (un)desired feature in the output. We assume contrastive prompt pairs encoding this feature:
\begin{equation}\label{eq:contrastive_dataset}
\Dpos := \{\pposn \mid n \in [\paircnt]\},\qquad \Dneg := \{\pnegn \mid n \in [\paircnt]\},
\end{equation}
where $\pposn$ contains the target feature and $\pnegn$ does not. 
\looseness-1Following activation steering \cite{turner2024activation, ode_activation_steering, skifstad2026local}, we assume white-box access to T2V weights and activations. Specifically, for all intervention layers and timesteps, we can access $x_{l,t}$ for $l\in[L]$ and $t\in[T]$, observe these during inference, and apply additive corrections. Our goal is to find a time-varying control policy $\pi:[L]\times[T]\times\mathcal{X}\rightarrow\mathcal{U}$ mapping activations $x_{l,t}$ to corrections $u_{l,t} \in \mathcal{U}$ that induce the desired feature-level behavior in the video.

\vspace{-8pt}
\section{Methods}
\label{sec:methods}
\vspace{-8pt}

\looseness-1 We overview our method (Fig.~\ref{fig:overview}). We formulate T2V generation as a dynamical system (Sec.~\ref{sec:methods_dynamics}), enabling activation steering via online error feedback and optimal control with minimal perturbations. To avoid raw-space control, we identify a task-relevant latent space from contrastive vectors and model dynamics there (Sec.~\ref{subsec:latentactivationdynamics}). We then design a latent LQR steering controller (Sec.~\ref{subsec:latent_activation_lqr}) guided by a feature setpoint (Sec.~\ref{subsec:linearfeaturesetpoint}) and bound projection-induced performance loss (Sec.~\ref{sec:theory}).

\subsection{Formalizing T2V Generation as a Dynamical System}\label{sec:methods_dynamics}
\vspace{-5pt}

\looseness-1 We recast T2V inference as a dynamical system whose states are internal
DiT activations, enabling control-based activation interventions. In
\eqref{eq:transformer}, generation proceeds over diffusion timesteps
$t\in[T]$, each with $L$ transformer blocks and one scheduler, inducing activations
$\{x_{l,t}:t\in[T],\,l\in[L+1]\}$, where $x_{l,t}$ enters block $l$ at
timestep $t$. We flatten this grid via $s:=(t-1)(L+1)+l$ and set
$x_s:=x_{l,t}$. The trajectory has length $T(L+1)$, with transitions
from transformer blocks and sampler updates. For intermediate layers
within timestep $t$, the dynamics are: 
\begin{equation}\label{eq:dynamics_transformer}
    x_{s+1}
    =
    f_s(x_s)
    :=
    \phi^{(l)}(x_{l,t}),
    \qquad l < L.
\end{equation}
At the end of a denoising pass, i.e., after block $L$, the DiT output is unpatchified and passed to the scheduler to produce the next latent video state $\hat x_{t-1}$, which is patchified and set as the next state: 
\begin{equation}\label{eq:dynamics_transition}
    x_{s+1}
    =
    x_{t-1, 1}
    =
    \texttt{Patchify}(\hat x_{t-1}); \qquad \hat x_{t-1}
    =
    \texttt{Scheduler}
    \left(
        \hat x_t,
        M(\hat x_t,h,t),
        t
    \right)
\end{equation}
Thus, the full T2V sampling process can be viewed as a nonlinear, time-varying dynamical system
\begin{equation}\label{eq:dynamics_total}
    x_{s+1} = f_s(x_s),
    \qquad s = 0,\ldots,T(L+1)-1,
\end{equation}
where $f_s$ alternates between transformer blocks \eqref{eq:dynamics_transformer} and scheduler updates \eqref{eq:dynamics_transition}.
We steer by adding control inputs in \eqref{eq:dynamics_total}. We study two intervention types: perturbing 1) the video-token activations after a transformer block, or 2) the text-embeddings used by cross-attention. Let $u_s^v$ and $u_s^h$ denote perturbations to the video-token and text-embedding streams, respectively. The steered transformer block is
\begin{subequations}\label{eq:steered_transformer}
\begin{align}
    \hspace{-8pt}M_{\mathrm{steer}}
    &:=
    \texttt{Unpatchify}
    \circ
    \rho^{(L)}
    \circ \cdots \circ
    \rho^{(1)}
    \circ
    \texttt{Patchify}(\hat x_t),
    \label{eq:steered_transformer_M}
    \\
    \hspace{-8pt}x_{l+1,t}%^{(k+1)}
    &:=
    \rho^{(l)}(x_{l,t},u_s^h,u_s^v)
    :=
    x_{\mathrm{attn}}
    +
    \texttt{MLP}_l
    \bigl(
        \texttt{Norm}(x_{\mathrm{attn}})
    \bigr)
    +
    u_s^v, \quad x_{1,t}:= \texttt{Patchify}(\hat x_t),
    \label{eq:steered_transformer_dynamics}
    \\
    \hspace{-8pt}x_{\mathrm{attn}}
    &:=
    x'
    +
    \texttt{CrossAttn}_l
    \bigl(
        \texttt{Norm}(x'),
        % h + u_s^h
        h_s + u_s^h
    \bigr),
    \quad\
    x'
    :=
    x_{l,t}
    +
    \texttt{SelfAttn}_l
    \bigl(
        \texttt{Norm}(x_{l,t})
    \bigr),
    \label{eq:steered_transformer_z}
\end{align}
\end{subequations}
where $h_s = h_{s-1}+u_{s-1}^h$. The combined control input is $u_s := (u_s^h,u_s^v) \in \mathcal{U}$, and $u_s^v$ may be applied to all video tokens, a selected subset of frames, a selected subset of spatial patches, or a selected set of layers and timesteps. Similarly, $u_s^h$ may be applied to all text tokens or only to tokens associated with the feature being steered. In our main experiments, we steer only the text embeddings at all layers and timesteps for steering efficacy, though App. \ref{app:multi_steer} provides examples with combined video-token and text-embedding steering.
Under these interventions, the controlled dynamics are
\begin{equation}\label{eq:controlled_dynamics}
    x_{s+1}
    =
    f_s(x_s,u_s).
\end{equation}


\vspace{-8pt}
\subsection{Constructing Reduced-Order Latent Activation Dynamics}
\label{subsec:latentactivationdynamics}
\vspace{-5pt}

From Sec.~\ref{sec:methods_dynamics}, T2V activation steering can be viewed as a control problem: at each flattened layer-timestep index $s$, the current activation is observed, a correction $u_s$ is applied, and the model evolves according to the controlled dynamics. LQR provides a reliable model-based method to find such steering policies. However, directly applying full activation-space LQR, as in \cite{skifstad2026local}, is computationally infeasible for T2V models because their activation states are orders of magnitude larger than LLMs'.

\looseness-1For example, Wan 2.1-14B~\cite{wan2025wan} generating $\pixframecnt=41$ frames at 480p has $\pixh=480$ and $\pixw=832$. After spatiotemporal compression, the DiT activation has shape $(\latframecnt,\lath,\latw,\hiddendim)=(11,30,52,5120)$, so 
$\actdim = 11\cdot 30\cdot 52\cdot 5120 = 87{,}859{,}200$. 
A full-space LQR Jacobian would have size $\actdim\times\actdim$, requiring over 30PB in single precision, making direct storage or manipulation infeasible. 
We address this by assuming the target feature lies mainly in a \textit{low-dimensional subspace}, with $\latdim\ll\actdim$, and validate this in Sec.~\ref{sec:results}. To construct the subspace, we collect vectorized activations $x_{l,t}^{(p)}\in\mathbb{R}^{\actdim}$ at each layer $l$ and timestep $t$ for prompts $p\in\Dpos\cup\Dneg$. The paired prompts $(\pposn,\pnegn)$ differ primarily by the target feature, so their activation differences isolate feature-relevant directions. 
Then, we construct the contrastive activation matrix $C_{l,t}\in\R^{\paircnt\times\actdim}$, whose $n$th row is the paired raw contrastive vector $C_{l,t}[n,:] = x_{l,t}^{(\pposn)}-x_{l,t}^{(\pnegn)}$. 
Since $\actdim$ is too large to form dense Jacobians, we use randomized SVD~\cite{halko2011finding} to obtain a low-rank contrastive basis. We partition the $L\times T$ layer-timestep indices into $R$ disjoint groups $\{\mathcal{I}_r\}_{r=1}^R$, each sharing one basis. For each $r\in[R]$, we run streaming randomized SVD over prompt pairs $n$ and 
$(l,t)\in\mathcal{I}_r$ on the GPU, accumulating sketches of the contrastive activation matrix, yielding an estimated rank-$\latdim$ right singular subspace without materializing any $\actdim\times\actdim$ matrices. The result is an orthonormal basis $ V_r \in \mathbb{R}^{\actdim\times \latdim}$ spanning the dominant contrastive directions in partition $r$. For each $(l,t)$, we compute the mean raw contrastive activation vector 
\begin{equation}\label{eq:mean_contrastive}
    \mu_{l,t}
    =
    \textstyle\frac{1}{\paircnt}
    \sum_{n=1}^{\paircnt}
    (x_{l,t}^{(\pposn)}-x_{l,t}^{(\pnegn)})
    % \big(x_{\pposn,l,t}-x_{\pnegn,l,t}\big)
    \in \R^{\actdim}.%\qquad \mu_k = \frac{1}{\paircnt} \sum_{n=1}^\paircnt (x_{\pposn,k}-x_{\pnegn,k})
\end{equation}
Let $r(l,t)$ denote the partition containing $(l,t)$ and set
$P_{l,t}:=V_{r(l,t)}^\top\in\R^{\latdim\times\actdim}$. Projecting the
mean contrastive vector gives
$c_{l,t}:=P_{l,t}\mu_{l,t}\in\R^{\latdim}$, the target feature direction
within the feature-relevant latent subspace $V_{r(l,t)}$. The captured
mean-contrast energy fraction is
\begin{equation}\label{eq:energy}
    \textstyle\rho_{l,t}
    =
    \frac{\|P_{l,t}\mu_{l,t}\|_2^2}{\|\mu_{l,t}\|_2^2}
    =
    \frac{\|V_{r(l,t)}^\top \mu_{l,t}\|_2^2}{\|\mu_{l,t}\|_2^2},
\end{equation}
serves as a diagnostic for whether the contrastive feature is well represented by the latent subspace.

For steering, we consider the T2V generation dynamics not on raw activations, but in the latent activation subspace. 
For each flattened layer-timestep index $s$, corresponding to a pair $(l,t)$ in inference order, we define the latent activation $z_s := P_s x_s \in \mathbb{R}^{\latdim}$, where $P_s := P_{l,t}=V_{r(l,t)}^\top \in \mathbb{R}^{\latdim\times\actdim}$, where $x_s:=x_{l,t}\in\mathbb{R}^{\actdim}$ is the vectorized raw activation and $P_s$ projects onto the rank-$\latdim$ contrastive subspace shared by the partition containing $(l,t)$. To obtain dynamics amenable to LQR, we linearize the projected controlled dynamics locally in latent space. Let $A_s$ and $B_s$ denote the raw Jacobians along a nominal trajectory $\{\bar x_s,\bar u_s\}$ of the raw controlled dynamics \eqref{eq:controlled_dynamics}:
    $A_s :=\frac{\partial f_s}{\partial x}
    \big|_{\bar x_s,\bar u_s}$, $B_s :=
    \frac{\partial f_s}{\partial u}
    \big|_{\bar x_s,\bar u_s}.$
Since $z_{s+1}=P_{s+1}x_{s+1}$, perturbations around the latent nominal trajectory satisfy $\delta z_{s+1}
    \approx
    \Atilde_s \delta z_s
    +
    \Btilde_s \delta u_s$, with
    $\delta z_s := z_s-\bar z_s$ and 
    $\delta u_s := u_s-\bar u_s,$ 
where
\begin{equation}\textstyle
    \Atilde_s
    :=
    \frac{\partial z_{s+1}}{\partial z_s}
    =
    P_{s+1}A_sP_s^\top
    \in \mathbb{R}^{\latdim\times\latdim},
    \qquad
    \Btilde_s
    :=
    \frac{\partial z_{s+1}}{\partial u_s}
    =
    P_{s+1}B_s
    \in \mathbb{R}^{\latdim\times\textdim}.
\end{equation}
Note that any raw activation $x_s$ can be decomposed as
    $x_s = x_s^\perp + P_s^\top z_s$, with
    $x_s^\perp \in \Null(P_s)$ and 
    $P_s^\top z_s \in \Row(P_s)$,
and the reduced dynamics model tracks only the component in $\Row(P_s)$. Notably, neither $A_s$ nor $B_s$ is materialized explicitly. Instead, products with $\Atilde_s$ and $\Btilde_s$ are computed using Jacobian-vector products (JVP) or vector-Jacobian products (VJP) through the DiT, followed by projection with $P_{s+1}$. For text-embedding control, we assume for simplicity that the same intervention $u_s\in\mathbb{R}^{\textdim}$ is added to every text token embedding, where $\textdim$ is the dimension of a single text-token embedding, which is usually on the order of a few thousand and is thus computationally feasible. This keeps $\Btilde_s$ small, with size $\latdim\times\textdim$. The resulting reduced-order LTV approximation is
\begin{equation}
\label{eq:latentdynamics}
    \delta z_{s+1}
    \approx
    \Atilde_s \delta z_s
    +
    \Btilde_s \delta u_s,
    \qquad
    s=0,\ldots,T(L+1)-1.
\end{equation}


\vspace{-10pt}
\subsection{Latent Linear Feature Setpoints (LLFS)}
\label{subsec:linearfeaturesetpoint}
\vspace{-6pt}

Within the latent space, we can define a set of \textit{latent linear feature setpoints} (LLFS) which define a desired feature strength level.
Let $e_s^z$ (resp. $v_s^z$)
\begin{equation}\label{eq:projected_contrastive}
    e^z_s := \textstyle\frac{1}{N}\sum_{n=1}^N P_s\big(x^{(p_n^+)}_s - x^{(p_n^-)}_s\big)
       = P_s \mu_s,
\quad
v^z_s := \frac{e^z_s}{\|e^z_s\|_2}.
\end{equation}
be the feature direction (resp. normalized feature direction) at the latent activations $z_s$, acquired by projecting each row of $C_{l,t}$ to the latent space and taking the average over the $\paircnt$ prompt pairs.
Following \cite{skifstad2026local}, we define latent feature strength $\beta_s$ and latent linear feature setpoint (LLFS) $\beta_s^\ast$ as
\begin{equation}\label{eq:lfs}
    \beta_{s}^z:= {(v_{s}^z)}^\top z_s,\quad
    \beta_s^{z,\ast}:=\lambda \|e_s^z\|_2
\end{equation}
where $\lambda$ is the hyperparameter that determines how strongly we would like to enforce the feature. 

\vspace{-5pt}
\subsection{Latent Activation LQR for Reaching Setpoints}
\label{subsec:latent_activation_lqr}
\vspace{-5pt}

Using the LTV approximation in \eqref{eq:latentdynamics}, we synthesize
a controller that steers latent activations toward the LLFS setpoints in
\eqref{eq:lfs}. We linearize around a representative trajectory
$\{(\bar z_s,\bar u_s)\}_{s=1}^{\layercnt\cdot\timestepcnt}$ to obtain
$\{\Atilde_s,\Btilde_s\}$ for each layer-timestep index $s$. As
validated in Sec.~\ref{sec:results}, this approximation generalizes
across prompts, allowing the same LTV model and LQR gains to be reused
for novel prompts without online Jacobian recomputation. For a realized
latent activation $z_s$, define the LLFS tracking error 
\begin{equation}
    \alpha_s := \beta_s^{z,*} - (v_s^z)^\top z_s,
    \quad
    z'_s := z_s + \alpha_s v_s^z,
    \quad
    \delta z_s := z_s - z'_s = -\alpha_s v_s^z.
    \label{eq:lfs_delta}
\end{equation}
Thus, $z'_s$ is the minimum-norm latent perturbation achieving the
desired feature strength $\beta_s^\ast$. We solve the LQR tracking
problem \eqref{eq:lqr_tracking} under the latent dynamics
\eqref{eq:latentdynamics}. 
The Riccati recursion yields feedback gains
$K_s\in\R^{\textdim\times\latdim}$, where $\textdim$ is the
dimension of the intervention $u_s$. The resulting Latent Activation-LQR (LA-LQR)
controller adapts to the LLFS tracking error \eqref{eq:lfs_delta} and can be written as:
\begin{equation}\label{eq:alqr}
    u_s^\ast
    :=
    \bar u_s - K_s\delta z_s
    =
    \bar u_s + K_s \alpha_s v_s 
    =
    \bar u_s + (\beta_s^{z,\ast} - {v_s^z}^\top z_s)K_sv_s^z .
\end{equation}
When $\bar u_s=0$, the intervention magnitude is proportional to the
online LLFS error $\beta_s^\ast-v_s^\top z_s$. Unlike open-loop
contrastive addition, it \textit{adapts online} to the realized latent activation at each
layer-timestep index. Since LQR is solved in latent space, Riccati
recursion scales with $\latdim\ll\actdim$, improving efficiency. After precomputing gains $K_s$,
runtime steering only requires projecting activations $z_s=P_sx_s$ and applying the matrix operations in 
\eqref{eq:alqr}. More generally, the formulation tracks any latent
reference $\{z_s^\ast\}$ by setting $\delta z_s=z_s-z_s^\ast$; LLFS uses
$z_s^\ast=z_s+(\beta_s^\ast-v_s^\top z_s)v_s$.

\subsection{Theoretical Analysis of LA-LQR}
\label{sec:theory}

We now relate a desired linear feature setpoint in the raw activation
space to an equivalent setpoint in the latent activation space.
Let $x_s\in\R^{\actdim}$ denote the raw activation at flattened
layer-timestep index $s$, and let
$z_s = P_s x_s$, where $P_s\in\R^{\latdim\times \actdim}$,
where $P_sP_s^\top=I_{\latdim}$. Let the orthogonal projection
onto the retained subspace be denoted $\Pi_s := P_s^\top P_s$.
Let $v^x_s := \frac{\mu_s}{\Vert \mu_s \Vert_2}\in\R^{\actdim}$ be a unit-normalized raw contrastive mean direction, where
$\mu_s$ is from \eqref{eq:mean_contrastive}. The raw feature strength of some activation $x_s$ is thus
$\beta^x_s := (v^x_s)^\top x_s$. 
Assume that the projection $P_s$ retains a nonzero component of this
feature direction $v^x_s$, i.e., $\gamma_s := \|P_s v^x_s\|_2 = \|\Pi_s v^x_s\|_2 = \sqrt{\rho_s} >0$, where $\rho_s$ is  the mean-contrast energy defined in \eqref{eq:energy}.
We then define the corresponding latent feature direction following \eqref{eq:projected_contrastive}, i.e.,
    $v^z_s := \frac{P_s v^x_s}{\|P_s v^x_s\|_2}
    =
    \frac{P_s v^x_s}{\gamma_s}$,
and the latent feature strength as $\beta^z_s := (v^z_s)^\top z_s$, following \eqref{eq:lfs}.

\begin{lemma}[Projection-calibrated setpoints]
\label{lem:projection_calibrated_lfs}
\looseness-1For any raw $x_s$, the raw and latent feature strengths
satisfy $\beta^x_s = \gamma_s \beta^z_s + \eta_s(x_s)$, 
where $\eta_s(x_s) := \bigl((I-\Pi_s)v^x_s\bigr)^\top x_s$
is the feature component lost by the projection. Thus, for
any desired raw-space setpoint $\beta^{x,\star}_s$, any scalar
$\bar\eta_s$ defines a latent-space setpoint
\begin{equation}
    \beta^{z,\star}_s
    :=
    (\beta^{x,\star}_s-\bar\eta_s)/\gamma_s = (\beta^{x,\star}_s-\bar\eta_s)/\sqrt{\rho_s}
    \label{eq:beta_z_from_beta_x}
\end{equation}
such that $\left|
    \beta^x_s-\beta^{x,\star}_s
    \right|
    \le
    \gamma_s
    \left|
    \beta^z_s-\beta^{z,\star}_s
    \right|
    +
    \left|
    \eta_s(x_s)-\bar\eta_s
    \right|$.
\end{lemma}


Lemma~\ref{lem:projection_calibrated_lfs} states that in order to track a raw-space threshold
$\beta^{x,\star}_s$, it suffices to track the latent threshold
$\beta^{z,\star}_s$ in \eqref{eq:beta_z_from_beta_x}. The only
irreducible discrepancy is the null-space term
$\eta_s(x_s)-\bar\eta_s$, which captures how much information from the raw
feature direction is discarded by the projection. In practice,
$\bar\eta_s$ can be set to zero or estimated from a calibration
trajectory, e.g.,
    $\bar\eta_s
    =
    \frac{1}{|\mathcal{D}_{\mathrm{cal}}|}
    \sum_{p\in\mathcal{D}_{\mathrm{cal}}}
    \eta_s(x_s(p))$.

We next show that LQR tracking in the latent space approximates raw
feature-setpoint tracking, up to linearization error and projection
loss. Let the exact projected dynamics be written as
\begin{equation}
    z_{s+1}
    =
    g_s(z_s,u_s)+\omega_s, \qquad g_s(z_s,u_s)
    :=
    P_{s+1} f_s(P_s^\top z_s+\bar n_s,u_s),
    \label{eq:projected_exact_dynamics}
\end{equation}
where $g_s(z_s,u_s)$ lifts $z_s$ using a nominal null-space component
$\bar n_s\in\Null(P_s)$, while $\omega_s$ captures unmodeled
null-space effects, with $\|\omega_s\|_2\le \xi_s$. Let
$\{(\bar z_s,\bar u_s)\}$ satisfy
$\bar z_{s+1}=g_s(\bar z_s,\bar u_s)$, and let
$\{\tilde A_s,\tilde B_s\}$ denote Jacobian linearizations about this
trajectory.
For perturbations $\delta z_s:=z_s-\bar z_s$ and
$\delta u_s:=u_s-\bar u_s$, define the Taylor remainder 
    $r_s(\delta z_s,\delta u_s)
    :=
    g_s(\bar z_s+\delta z_s,\bar u_s+\delta u_s)
    -
    g_s(\bar z_s,\bar u_s)
    -
    \tilde A_s\delta z_s
    -
    \tilde B_s\delta u_s$.
Further assume that, in a neighborhood of the nominal trajectory, there exists a Lipschitz constant $L_s \ge 0$ such that
    $\|r_s(\delta z_s,\delta u_s)\|_2
    \le
    \frac{L_s}{2}
    \big\|
    \begin{bmatrix}
    \delta z_s,\
    \delta u_s
    \end{bmatrix}^\top
    \big\|_2^2$.

\begin{theorem}[Latent closed-loop tracking under projection loss]
\label{thm:latent_tracking_projection_loss}
Consider the projected dynamics \eqref{eq:projected_exact_dynamics}.
Apply the LA-LQR feedback law
$\delta u_s=-K_s\delta z_s$,
and define
$\widehat A_s:=\tilde A_s-\tilde B_sK_s$.
Let $\widehat\Phi_{s,j} := \widehat A_{s-1}\widehat A_{s-2}\cdots \widehat A_j,$ if $s>j$, and $I$ if $s = j$. 
Then, for all $s$,
\begin{equation}
    \textstyle\|\delta z_s\|_2
    \le
    \|\widehat\Phi_{s,1}\|_2\|\delta z_1\|_2
    +
    \sum_{i=1}^{s-1}
    \|\widehat\Phi_{s,i+1}\|_2
    \big(
    \xi_i
    +
    \frac{L_i}{2}
    \big\|
    \begin{bmatrix}
    \delta z_i,\ \
    -K_i\delta z_i
    \end{bmatrix}^\top
    \big\|_2^2
    \big).
    \label{eq:latent_tracking_projection_bound_2}
\end{equation}
\end{theorem}


\begin{corollary}[Raw LFS tracking by LA-LQR]
\label{cor:raw_lfs_tracking_latent_lqr}
Let the assumptions of Lemma~\ref{lem:projection_calibrated_lfs} and
Theorem~\ref{thm:latent_tracking_projection_loss} hold. Suppose the
latent nominal trajectory is constructed to satisfy
$(v^z_s)^\top \bar z_s = \beta^{z,\star}_s$, 
where $\beta^{z,\star}_s$ is chosen according to
\eqref{eq:beta_z_from_beta_x}. Then the raw-space feature tracking
error $\epsilon^x_s
    :=
    (v^x_s)^\top x_s-\beta^{x,\star}_s$
satisfies
\[
\begin{aligned}
    |\epsilon^x_s|
    \le\;
    \gamma_s
    \left|
    (v^z_s)^\top \widehat\Phi_{s,1}\delta z_1
    \right|
    +
    \gamma_s
    \sum_{i=1}^{s-1}
    \left\|
    (v^z_s)^\top \widehat\Phi_{s,i+1}
    \right\|_2
    \left(
    \xi_i
    +
    \frac{L_i}{2}
    \left\|
    \begin{bmatrix}
    \delta z_i\\
    -K_i\delta z_i
    \end{bmatrix}
    \right\|_2^2
    \right)
    +
    \left|
    \eta_s(x_s)-\bar\eta_s
    \right|.
\end{aligned}
\label{eq:raw_lfs_tracking_bound}
\]
\end{corollary}

Corollary~\ref{cor:raw_lfs_tracking_latent_lqr} shows that LA-LQR approximates raw-space LFS tracking when the latent closed-loop
tracking error, projected dynamics error $\xi_i$, and discarded raw
feature component $\eta_s(x_s)-\bar\eta_s$ are small. If
$v^x_s\in\Row(P_s)$ and the projected dynamics are exact, then
$\gamma_s=1$, $\eta_s(x_s)=0$, and $\xi_i=0$, recovering the standard
latent feature-tracking bound \cite[Eq. 23]{skifstad2026local}.

\section{Results}\label{sec:results}
\vspace{-7pt}

We first evaluate the validity of the LTV approximation to the latent dynamics \eqref{eq:latentdynamics} and assess whether the latent space preserves the information contained in raw contrastive vectors from the activation space. We then evaluate LA-LQR's ability to steer new concepts into generated videos. Finally, we provide quantitative and qualitative evaluations of T2V safeguarding against harmful prompts.

\begin{table}[t]
\centering
\caption{Numerical evaluations on T2VSafetyBench \cite{miao2024t2vsafetybench} on the Wan2.1-T2V-14B model \cite{wan2025wan}.}
\label{tab:t2vsafetybench_vertical}
\resizebox{\textwidth}{!}{%
\begin{tabular}{ll ccccc}
\toprule
\textbf{Metric} & \textbf{Category} & \textbf{Wan} & \cite{yoon2024safree} & \cite{facchiano2025video} & \cite{ekin2026unreasonable} & \textbf{Ours} \\
\midrule

\multirow{6}{*}{\textbf{Violation Rate} $\downarrow$} 
 & Copyright \& Trademarks  & 71.0\% $\pm$ 0.032 & 51.5\% $\pm$ 0.035 & 71.0\% $\pm$ 0.032 & 65.0\% $\pm$ 0.034 & \textbf{37.0\% $\pm$ 0.034} \\
 & Pornography              & 50.0\% $\pm$ 0.035 & 36.5\% $\pm$ 0.034 & 50.0\% $\pm$ 0.035 & 47.5\% $\pm$ 0.035 & \textbf{9.5\% $\pm$ 0.021} \\
 & Gore                     & 42.0\% $\pm$ 0.035 & 27.5\% $\pm$ 0.032 & 39.5\% $\pm$ 0.035 & 31.0\% $\pm$ 0.033 & \textbf{14.0\% $\pm$ 0.025} \\
 & Public Figure            & 10.5\% $\pm$ 0.022 & 6.0\% $\pm$ 0.017  & 10.0\% $\pm$ 0.021 & 5.0\% $\pm$ 0.015  & \textbf{3.0\% $\pm$ 0.012} \\
 & Sequential Action Risk   & 10.9\% $\pm$ 0.042 & 9.1\% $\pm$ 0.039  & 10.9\% $\pm$ 0.042 & 10.9\% $\pm$ 0.042 & \textbf{7.3\% $\pm$ 0.035} \\
 \cmidrule{2-7}
 & \textbf{Average}         & 36.9\% $\pm$ 0.033 & 26.1\% $\pm$ 0.031 & 36.3\% $\pm$ 0.033 & 31.9\% $\pm$ 0.032 & \textbf{14.6\% $\pm$ 0.025} \\

\midrule
\multirow{6}{*}{\begin{tabular}[c]{@{}l@{}}\textbf{VBench (Subject} \\ \textbf{Consistency)} $\uparrow$\end{tabular}} 
 & Copyright \& Trademarks  & 0.977 $\pm$ 0.019 & 0.973 $\pm$ 0.021 & 0.976 $\pm$ 0.019 & 0.976 $\pm$ 0.021 & 0.976 $\pm$ 0.021 \\
 & Pornography              & 0.973 $\pm$ 0.025 & 0.973 $\pm$ 0.024 & 0.972 $\pm$ 0.029 & 0.972 $\pm$ 0.027 & 0.974 $\pm$ 0.025 \\
 & Gore                     & 0.946 $\pm$ 0.051 & 0.963 $\pm$ 0.035 & 0.948 $\pm$ 0.050 & 0.944 $\pm$ 0.052 & 0.974 $\pm$ 0.014 \\
 & Public Figure            & 0.972 $\pm$ 0.018 & 0.971 $\pm$ 0.018 & 0.972 $\pm$ 0.017 & 0.971 $\pm$ 0.019 & 0.970 $\pm$ 0.014 \\
 & Sequential Action Risk   & 0.949 $\pm$ 0.040 & 0.953 $\pm$ 0.039 & 0.948 $\pm$ 0.044 & 0.948 $\pm$ 0.041 & 0.962 $\pm$ 0.030 \\
 \cmidrule{2-7}
 & \textbf{Average}         & 0.964 $\pm$ 0.031 & 0.967 $\pm$ 0.027 & 0.963 $\pm$ 0.032 & 0.962 $\pm$ 0.032 & 0.971 $\pm$ 0.021 \\

\midrule
\multirow{6}{*}{\textbf{CAPS} $\uparrow$} 
 & Copyright \& Trademarks  & N/A & 0.672 $\pm$ 0.161 & 0.848 $\pm$ 0.058 & 0.804 $\pm$ 0.150 & 0.674 $\pm$ 0.127 \\
 & Pornography              & N/A & 0.619 $\pm$ 0.203 & 0.802 $\pm$ 0.152 & 0.797 $\pm$ 0.143 & 0.588 $\pm$ 0.194 \\
 & Gore                     & N/A & 0.691 $\pm$ 0.134 & 0.817 $\pm$ 0.053 & 0.780 $\pm$ 0.093 & 0.408 $\pm$ 0.116 \\
 & Public Figure            & N/A & 0.679 $\pm$ 0.137 & 0.829 $\pm$ 0.054 & 0.691 $\pm$ 0.144 & 0.660 $\pm$ 0.673 \\
 & Sequential Action Risk   & N/A & 0.693 $\pm$ 0.123 & 0.843 $\pm$ 0.068 & 0.854 $\pm$ 0.045 & 0.733 $\pm$ 0.100 \\
 \cmidrule{2-7}
 & \textbf{Average}         & N/A & 0.671 $\pm$ 0.152 & 0.828 $\pm$ 0.077 & 0.785 $\pm$ 0.115 & 0.6126 $\pm$ 0.242 \\
 
\bottomrule
\end{tabular}%
}   \vspace{-8pt}
\end{table}

\vspace{-10pt}
\begin{table}[t]
\centering
\caption{Numerical evaluations on SafeSora \cite{dai2024safesora} on the HunyuanVideo-1.5 model \cite{kong2024hunyuanvideo}.}
\label{tab:t2vsafetybench_hunyuan_vertical}
\resizebox{\textwidth}{!}{%
\begin{tabular}{ll ccccc}
\toprule
\textbf{Metric} & \textbf{Category} & \textbf{Hunyuan} & \cite{yoon2024safree} & \cite{facchiano2025video} & \cite{ekin2026unreasonable} & \textbf{Ours} \\
\midrule

\multirow{6}{*}{\textbf{Violation Rate} $\downarrow$} 
 & Violence               & 32.0\% $\pm$ 0.038 & 30.7\% $\pm$ 0.036 & 32.5\% $\pm$ 0.036 & 33.7\% $\pm$ 0.037 & \textbf{22.2\% $\pm$ 0.034} \\
 & Terrorism              & 32.0\% $\pm$ 0.093 & 24.0\% $\pm$ 0.085 & \textbf{16.0\% $\pm$ 0.073} & 24.0\% $\pm$ 0.085 & 24.0\% $\pm$ 0.085 \\
 & Racism                 & 12.5\% $\pm$ 0.068 & 13.3\% $\pm$ 0.051 & 24.4\% $\pm$ 0.064 & 26.7\% $\pm$ 0.066 & \textbf{0.0\% $\pm$ 0.000} \\
 & Sexual                 & 43.8\% $\pm$ 0.088 & 45.5\% $\pm$ 0.087 & 45.5\% $\pm$ 0.087 & 45.5\% $\pm$ 0.087 & \textbf{0.0\% $\pm$ 0.000} \\
 & Animal Abuse           & 37.0\% $\pm$ 0.093 & 29.6\% $\pm$ 0.088 & 44.4\% $\pm$ 0.096 & 33.3\% $\pm$ 0.091 & \textbf{14.8\% $\pm$ 0.068} \\
 \cmidrule{2-7}
 & \textbf{Average}       & 31.5\% $\pm$ 0.076 & 28.6\% $\pm$ 0.069 & 32.6\% $\pm$ 0.071 & 32.6\% $\pm$ 0.073 & \textbf{12.2\% $\pm$ 0.037} \\

\midrule
\multirow{6}{*}{\begin{tabular}[c]{@{}l@{}}\textbf{VBench (Subject} \\ \textbf{Consistency)} $\uparrow$\end{tabular}} 
 & Violence               & 0.942 $\pm$ 0.035 & 0.922 $\pm$ 0.038 & 0.929 $\pm$ 0.036 & 0.922 $\pm$ 0.039 & 0.971 $\pm$ 0.008 \\
 & Terrorism              & 0.925 $\pm$ 0.039 & 0.910 $\pm$ 0.048 & 0.917 $\pm$ 0.033 & 0.913 $\pm$ 0.046 & 0.959 $\pm$ 0.018 \\
 & Racism                 & 0.960 $\pm$ 0.021 & 0.931 $\pm$ 0.032 & 0.939 $\pm$ 0.027 & 0.937 $\pm$ 0.027 & 0.963 $\pm$ 0.014 \\
 & Sexual                 & 0.966 $\pm$ 0.026 & 0.933 $\pm$ 0.036 & 0.937 $\pm$ 0.040 & 0.935 $\pm$ 0.035 & 0.967 $\pm$ 0.016 \\
 & Animal Abuse           & 0.951 $\pm$ 0.025 & 0.917 $\pm$ 0.042 & 0.943 $\pm$ 0.039 & 0.932 $\pm$ 0.046 & 0.975 $\pm$ 0.011 \\
 \cmidrule{2-7}
 & \textbf{Average}       & 0.949 $\pm$ 0.029 & 0.923 $\pm$ 0.039 & 0.933 $\pm$ 0.035 & 0.928 $\pm$ 0.039 & 0.967 $\pm$ 0.013 \\

\midrule
\multirow{6}{*}{\textbf{CAPS} $\uparrow$} 
 & Violence               & N/A & 0.729 $\pm$ 0.080 & 0.740 $\pm$ 0.085 & 0.737 $\pm$ 0.091 & 0.582 $\pm$ 0.125 \\
 & Terrorism              & N/A & 0.719 $\pm$ 0.085 & 0.698 $\pm$ 0.085 & 0.719 $\pm$ 0.068 & 0.696 $\pm$ 0.088 \\
 & Racism                 & N/A & 0.714 $\pm$ 0.104 & 0.741 $\pm$ 0.086 & 0.752 $\pm$ 0.090 & 0.706 $\pm$ 0.107 \\
 & Sexual                 & N/A & 0.665 $\pm$ 0.239 & 0.744 $\pm$ 0.148 & 0.677 $\pm$ 0.224 & 0.546 $\pm$ 0.221 \\
 & Animal Abuse           & N/A & 0.730 $\pm$ 0.079 & 0.771 $\pm$ 0.053 & 0.742 $\pm$ 0.079 & 0.694 $\pm$ 0.111 \\
 \cmidrule{2-7}
 & \textbf{Average}       & N/A & 0.711 $\pm$ 0.117 & 0.739 $\pm$ 0.091 & 0.725 $\pm$ 0.110 & 0.645 $\pm$ 0.130 \\
 
\bottomrule
\end{tabular}%
}\vspace{-10pt}
\end{table}


\paragraph{Validation of Linearity and Compressed Feature Fidelity}


\begin{figure}
    \centering\vspace{-20pt}
    \includegraphics[width=\linewidth]{figs/energy_norm_combined.pdf}\vspace{-27pt}
    \caption{\textbf{(a)} Proportion of energy of $C_r$ matrices captured in the subspace spanned by the top-$D_\mathrm{lat}=64$ right singular vectors, over $(t,l)$, for pornography feature. \textbf{(b)} Normalized Frobenius norm between \textit{(Top)} $A_s$ computed from 20 different prompts and \textit{(Bottom)} random matrices, at \textit{(left)} layer 5, \textit{(middle)} layer 25, and \textit{(right)} layer 35.\vspace{-10pt}}
    \label{fig:energy}
\end{figure}



\looseness-1We validate the assumption that features lie in a low-dimensional subspace by inspecting the energy of each $C_{(l, t)}$.
Fig. \ref{fig:energy} (a) shows that a large portion of energy is captured by the first $D_{\mathrm{lat}}=64$ singular vectors across most of $(t,l)$.
In (b), we visualize the normalized Frobenius norm of the difference between Jacobians $A_{(l,t)}^{(i)}$ and $A_{(l,t)}^{(j)}$, defined as $\|A_{(l,t)}^{(i)}-A_{(l,t)}^{(j)}\|_F/\left(\frac{1}{M}\sum_m\|A_{(l,t)}^{(m)}\|_F\right)$, for each unique pair $(i, j)$ of prompts from a set of $M=20$ prompts.
Compared to those computed with random matrices, the Frobenius norm is significantly smaller, which demonstrates the similarity of Jacobians across varying prompts, motivating the reuse of a Jacobian matrix computed on one prompt to control the generation for another.

\vspace{-7pt}
\paragraph{Concept Steering}

We demonstrate concept steering on the \texttt{Wan2.1-T2V-14B lightx2v 4-step-distill LoRA} \cite{wan2025wan} model by steering toward the concept of the color red. We construct 20 contrastive prompts by adding the word ``red" to selected nouns in each neutral prompt. Fig. \ref{fig:concept_steering} shows the first frame of videos generated by LA-LQR under different LQR parameters that control steering strength, starting from a cyberpunk-themed prompt (see App. \ref{app:concept} for details). Increasing the LQR state-cost parameter $Q$ produces stronger steering. As $Q$ increases, red becomes progressively more prominent, demonstrating effective concept-level control. See App. \ref{app:multi_steer} for examples of steering video tokens alone, as well as steering both video and text tokens jointly (Figs. \ref{fig:video_steer_only} and \ref{fig:text_video_steer}). 

\begin{figure}
    \centering
    \includegraphics[width=\linewidth]{figs/fig_concept.pdf}\vspace{-20pt}
    \caption{Steering the concept of ``the color red" in Wan, with increasing steering strength. \textbf{(a)} Unsteered. \textbf{(b)} $Q = I$. \textbf{(c)} $Q = 1.5I$. \textbf{(d)} $Q = 2I$. \textbf{(e)} $Q = 2.5I$. \textbf{(f)} $Q = 3I$.\vspace{-13pt} }
    \label{fig:concept_steering}
\end{figure}

\vspace{-7pt}
\paragraph{Safeguarding T2V Models}

\begin{figure}
    \centering
    \includegraphics[width=\linewidth]{figs/fig_wan_2x2.pdf}\vspace{-20pt}
    \caption{\looseness-1Safeguarding Wan against harmful prompts from T2VSafetyBench~\cite{miao2024t2vsafetybench}. For each example, the \textit{left column} shows the unsteered generation and the \textit{right column} shows the LA-LQR-steered generation. The \textit{top, middle, and bottom rows} show the first, middle, and final frames, respectively. We consider four categories: \textbf{(a)} pornography, \textbf{(b)} gore, \textbf{(c)} public figures, and \textbf{(d)} sequential action risk.\vspace{-15pt}
    }
    \label{fig:wan_steering}
\end{figure}

We evaluate our method on safeguarding T2V models. We evaluate on two state-of-the-art T2V models, \texttt{Wan2.1-T2V-14B + lightx2v 4-step-distill LoRA} \cite{wan2025wan} and \texttt{HunyuanVideo-1.5 + lightx2v 4step LoRA} \cite{kong2024hunyuanvideo}.

\noindent \textit{Metrics:\quad} We evaluate using three quantitative metrics. \textit{(1) Violation Rate}. Following prior T2V safety work \cite{miao2024t2vsafetybench,dai2024safesora}, we measure the percentage of generations containing unsafe content across predefined categories; \textbf{this is our primary metric}. We use a GPT-4o evaluator on each 41-frame video, following \cite[Fig.~3]{miao2024t2vsafetybench}, and aggregate to a binary per-video label. This evaluation aligns closely with human judgment \cite{miao2024t2vsafetybench}. In Tables~\ref{tab:t2vsafetybench_vertical} and~\ref{tab:t2vsafetybench_hunyuan_vertical}, \textit{Violation Rate} denotes the percentage of videos containing the undesired category. \textit{(2) Video Quality (VBench \cite{huang2024vbench})}. We use VBench \cite{huang2024vbench} to measure video quality and realism, including subject consistency, motion smoothness, and aesthetic quality. We report \textit{subject consistency} in Tables~\ref{tab:t2vsafetybench_vertical} and~\ref{tab:t2vsafetybench_hunyuan_vertical}, with full per-category metrics in App.~\ref{app:vbench}.
\textit{(3) Content Alignment Preservation Score (CAPS).} To measure semantic similarity between the steered video and its original unsteered counterpart, we use GPT-4o to generate semantic descriptions of both videos, then compute their relatedness using text embeddings from ChatGPT \texttt{text-embedding-3-small}.

\noindent \textit{Datasets:\quad} We evaluate on T2VSafetyBench \cite{miao2024t2vsafetybench} and SafeSora \cite{dai2024safesora}, which contain harmful prompts in categories such as pornography, gore, and copyright. We select five categories from each dataset.

\noindent \textit{Baselines:\quad} We compare against three T2V safeguarding/content-steering methods: (A) \cite{facchiano2025video}, which updates model weights using contrastive-vector updates resembling activation addition \cite{turner2024activation}; (B) SAFREE \cite{yoon2024safree}, which removes toxic directions from text embeddings during denoising; and (C) \cite{ekin2026unreasonable}, which uses an LLM to select text-token positions for steering in text-embedding space.

\noindent \textit{Qualitative Results:\quad} Across both benchmarks and models, LA-LQR reduces unsafe concepts. In Fig.~\ref{fig:wan_steering}, it steers Wan away from unsafe generations while often preserving prompt semantics. For instance, Fig.~\ref{fig:wan_steering}a shows a minimally invasive shift from graphic nudity to minimal clothing. Similar interventions appear in Fig.~\ref{fig:wan_steering}c, where Angela Merkel is replaced by a visually similar person, and Fig.~\ref{fig:wan_steering}d, where the man's position is made less precarious while preserving the background and color palette. Fig.~\ref{fig:wan_steering}b requires a larger intervention but preserves key scene elements, including a house, doors, and windows. Additional Hunyuan examples on SafeSora~\cite{dai2024safesora} appear in App.~\ref{app:qualitative_results}.

\noindent \textit{Quantitative Results:\quad} Table~\ref{tab:t2vsafetybench_vertical} shows that LA-LQR achieves the lowest violation rate on T2VSafetyBench. LA-LQR can steer aggressively toward a desired feature setpoint while applying only the minimum-norm perturbation needed to reach it. In contrast, \cite{facchiano2025video} must be tuned more conservatively: increasing its steering strength often introduces severe visual artifacts, e.g., Fig.~\ref{fig:oversteer}. Filtering- and embedding-based methods such as SAFREE~\cite{yoon2024safree} and \cite{ekin2026unreasonable} also struggle on T2VSafetyBench, where jailbroken prompts make input-level filtering and token selection unreliable. LA-LQR's CAPS scores are sometimes lower because removing harmful concepts can require larger semantic changes; however, these changes reflect successful safety interventions rather than degraded alignment. Similar trends hold for Hunyuan in Table~\ref{tab:t2vsafetybench_hunyuan_vertical}: LA-LQR reduces unsafe generations while maintaining strong VBench quality, indicating an effective safety-quality tradeoff across models.

\vspace{-8pt}
\section{Discussion, Limitations, and Conclusion}\label{sec:conclusion}
\vspace{-8pt}

LA-LQR is a reduced-order optimal-control framework for T2V steering that reduces unsafe concepts while preserving video quality and prompt semantics. By projecting activations into latent subspaces and applying closed-loop LQR feedback, it makes activation control tractable for modern video DiTs.

\vspace{-8pt}
\paragraph{Limitations.}
First, LA-LQR incurs memory and storage overhead: although it avoids full activation-space Jacobians, projection bases and projected dynamics must still be stored across layers and timesteps. Coarser basis sharing or quantized projections could reduce this cost. Second, performance depends on latent dimension and contrastive-subspace quality; if contrastive prompts miss the target concept, projection may discard relevant directions. Monitoring $\rho_{l,t}$ can guide adaptive rank increases or additional contrastive-pair collection. Third, LA-LQR introduces steering hyperparameters $Q$, $R$, $Q_H$, and $\lambda$, whose effects are intuitive: larger $Q$ or $\lambda$, or smaller $R$, strengthens steering. Future work could automate tuning with \textit{adaptive} LQR controllers~\cite{abbasi2011regret} driven by detector feedback. Fourth, LA-LQR relies on local linear latent dynamics, which may degrade for adversarially-selected prompts; storing multiple local models may help. Finally, LA-LQR assumes white-box activation access and only guarantees embeddings satisfying the disturbance bound $\xi_s$ in Cor.~\ref{cor:raw_lfs_tracking_latent_lqr}. Thus, it should be paired with prompt filtering or output moderation; black-box surrogate extensions remain future work.

\vspace{-8pt}
\paragraph{Conclusion.}
We introduced LA-LQR, a reduced-order linear optimal-control method for activation steering in T2V models. LA-LQR treats inference as a dynamical system, projects activations into a task-relevant latent subspace, and computes timestep- and layer-specific feedback interventions toward desired feature setpoints while penalizing unnecessary perturbations. Across concept steering and T2V safeguarding, LA-LQR reduces unsafe generations relative to prior baselines while maintaining strong video quality and semantic preservation. More broadly, our results suggest control-theoretic activation steering as a promising direction for reliable, training-free alignment of T2V models.

\bibliographystyle{IEEEtran}
\bibliography{references}

\clearpage


%%%%%%%%%%%%%%%%%%%%%%%%%%%%%%%%%%%%%%%%%%%%%%%%%%%%%%%%%%%%

\appendix

\section{Proofs}\label{app:proofs}

\begin{lemma}[Projection-calibrated feature setpoints]
\label{lem:projection_calibrated_lfs_app}
\looseness-1For any raw $x_s$, the raw and latent feature strengths
satisfy $\beta^x_s = \gamma_s \beta^z_s + \eta_s(x_s)$, 
where $\eta_s(x_s) := \bigl((I-\Pi_s)v^x_s\bigr)^\top x_s$
is the feature component lost by the projection. Thus, for
any desired raw-space setpoint $\beta^{x,\star}_s$, any scalar
$\bar\eta_s$ defines a latent-space setpoint
\begin{equation}
    \beta^{z,\star}_s
    :=
    (\beta^{x,\star}_s-\bar\eta_s)/\gamma_s = (\beta^{x,\star}_s-\bar\eta_s)/\sqrt{\rho_s}
    \label{eq:beta_z_from_beta_x}
\end{equation}
such that $\left|
    \beta^x_s-\beta^{x,\star}_s
    \right|
    \le
    \gamma_s
    \left|
    \beta^z_s-\beta^{z,\star}_s
    \right|
    +
    \left|
    \eta_s(x_s)-\bar\eta_s
    \right|$.
\end{lemma}

\begin{proof}
Since $P_sP_s^\top=I$, we have
\[
    \beta^z_s
    =
    (v^z_s)^\top P_sx_s
    =
    \frac{(P_sv^x_s)^\top P_sx_s}{\gamma_s}
    =
    \frac{(v^x_s)^\top \Pi_s x_s}{\gamma_s}.
\]
Therefore,
\[
    \gamma_s\beta^z_s
    =
    (v^x_s)^\top \Pi_s x_s.
\]
Decomposing $x_s=\Pi_sx_s+(I-\Pi_s)x_s$ gives
\[
    \beta^x_s
    =
    (v^x_s)^\top x_s
    =
    (v^x_s)^\top \Pi_sx_s
    +
    (v^x_s)^\top (I-\Pi_s)x_s
    =
    \gamma_s\beta^z_s+\eta_s(x_s).
\]
Substituting \eqref{eq:beta_z_from_beta_x} and applying the triangle
inequality yields the result.
\end{proof}


\begin{theorem}[Latent closed-loop tracking under projection loss]
\label{thm:app_latent_tracking_projection_loss}
Consider the projected dynamics \eqref{eq:projected_exact_dynamics}.
Apply the LA-LQR feedback law
$\delta u_s=-K_s\delta z_s$,
and define
$\widehat A_s:=\tilde A_s-\tilde B_sK_s$.
Let $\widehat\Phi_{s,j} := \widehat A_{s-1}\widehat A_{s-2}\cdots \widehat A_j,$ if $s>j$, and $I$ if $s = j$. 
Then, for all $s$,
\begin{equation}
    \textstyle\|\delta z_s\|_2
    \le
    \|\widehat\Phi_{s,1}\|_2\|\delta z_1\|_2
    +
    \sum_{i=1}^{s-1}
    \|\widehat\Phi_{s,i+1}\|_2
    \big(
    \xi_i
    +
    \frac{L_i}{2}
    \big\|
    \begin{bmatrix}
    \delta z_i,\ \
    -K_i\delta z_i
    \end{bmatrix}^\top
    \big\|_2^2
    \big).
    \label{eq:latent_tracking_projection_bound_2}
\end{equation}
\end{theorem}

\begin{proof}
The closed-loop deviation dynamics are
\[
\begin{aligned}
    \delta z_{s+1}
    &=
    z_{s+1}-\bar z_{s+1} \\
    &=
    g_s(\bar z_s+\delta z_s,\bar u_s+\delta u_s)
    +\omega_s
    -
    \left(g_s(\bar z_s,\bar u_s)\right)\\
    &=
    \tilde A_s\delta z_s+\tilde B_s\delta u_s
    +
    r_s(\delta z_s,\delta u_s)
    +
    \omega_s.
\end{aligned}
\]
Substituting $\delta u_s=-K_s\delta z_s$ gives
\[
    \delta z_{s+1}
    =
    \widehat A_s\delta z_s
    +
    r_s(\delta z_s,-K_s\delta z_s)
    +
    \omega_s.
\]
Unrolling this recursion and applying submultiplicativity and the
triangle inequality gives:
\begin{equation}
    \|\delta z_s\|_2
    \le
    \|\widehat\Phi_{s,1}\|_2\|\delta z_1\|_2
    +
    \sum_{i=1}^{s-1}
    \|\widehat\Phi_{s,i+1}\|_2
    \left(
    \|r_i(\delta z_i,-K_i\delta z_i)\|_2
    +
    \xi_i
    \right).
    % \label{eq:latent_tracking_projection_bound_1}
\end{equation}
Applying the assumed
quadratic remainder bound and $\|\omega_i\|_2\le \xi_i$ gives
\begin{equation}
    \|\delta z_s\|_2
    \le
    \|\widehat\Phi_{s,1}\|_2\|\delta z_1\|_2
    +
    \sum_{i=1}^{s-1}
    \|\widehat\Phi_{s,i+1}\|_2
    \left(
    \xi_i
    +
    \frac{L_i}{2}
    \left\|
    \begin{bmatrix}
    \delta z_i\\
    -K_i\delta z_i
    \end{bmatrix}
    \right\|_2^2
    \right).
    % \label{eq:latent_tracking_projection_bound_2}
\end{equation}

\end{proof}

\begin{corollary}[Raw LFS tracking by LA-LQR]
\label{cor:app_raw_lfs_tracking_latent_lqr}
Let the assumptions of Lemma~\ref{lem:projection_calibrated_lfs} and
Theorem~\ref{thm:latent_tracking_projection_loss} hold. Suppose the
latent nominal trajectory is constructed to satisfy
$(v^z_s)^\top \bar z_s = \beta^{z,\star}_s$, 
where $\beta^{z,\star}_s$ is chosen according to
\eqref{eq:beta_z_from_beta_x}. Then the raw-space feature tracking
error $\epsilon^x_s
    :=
    (v^x_s)^\top x_s-\beta^{x,\star}_s$
satisfies
\[
\begin{aligned}
    |\epsilon^x_s|
    \le\;
    \gamma_s
    \left|
    (v^z_s)^\top \widehat\Phi_{s,1}\delta z_1
    \right|
    +
    \gamma_s
    \sum_{i=1}^{s-1}
    \left\|
    (v^z_s)^\top \widehat\Phi_{s,i+1}
    \right\|_2
    \left(
    \xi_i
    +
    \frac{L_i}{2}
    \left\|
    \begin{bmatrix}
    \delta z_i\\
    -K_i\delta z_i
    \end{bmatrix}
    \right\|_2^2
    \right)
    +
    \left|
    \eta_s(x_s)-\bar\eta_s
    \right|.
\end{aligned}
\label{eq:raw_lfs_tracking_bound}
\]
\end{corollary}

\begin{proof}
By Lemma~\ref{lem:projection_calibrated_lfs},
\[
    |\epsilon^x_s|
    \le
    \gamma_s
    \left|
    (v^z_s)^\top z_s-\beta^{z,\star}_s
    \right|
    +
    |\eta_s(x_s)-\bar\eta_s|.
\]
Since $(v^z_s)^\top \bar z_s=\beta^{z,\star}_s$,
\[
    (v^z_s)^\top z_s-\beta^{z,\star}_s
    =
    (v^z_s)^\top \delta z_s.
\]
Using the unrolled closed-loop deviation dynamics from
Theorem~\ref{thm:latent_tracking_projection_loss} and applying the
triangle inequality yields \eqref{eq:raw_lfs_tracking_bound}.
\end{proof}



\section{Implementation Details}\label{app:implementation}

\subsection{Model and Inference Settings}
The parameters and configurations for the models that we used for our experiments are provided in \Cref{tab:model_inference_settings}. 
\begin{table}[h]
\centering
\caption{Model and inference settings for the two LightX2V pipelines.}
\label{tab:model_inference_settings}
\begin{tabular}{lll}
\hline
\textbf{Parameter} & \textbf{Wan 2.1 LightX2V} & \textbf{HunyuanVideo 1.5 LightX2V} \\
\hline
Base model & Wan2.1-T2V-14B & HunyuanVideo-1.5, 480p T2V \\
Distillation & 4-step LightX2V LoRA, rank 64 & 4-step LightX2V checkpoint \\
Transformer depth & 40 DiT blocks & 54 double blocks \\
Hidden dimension & 5120 & 2048 \\
Sampling steps & 4 & 4 \\
Guidance scale & 1.0 & 1.0 \\
Scheduler / flow shift & UniPC, shift 5.0 & Flow shift 9.0 \\
Resolution & $832 \times 480$ & $848 \times 480$ \\
Frames & 41 & 41 \\
FPS & -- & 16 \\
Random seed & 42 & 42 \\
Main dtype & bfloat16 & bfloat16 \\
Jacobian accumulation dtype & fp32 & fp32 \\
\hline
\end{tabular}
\end{table}

\subsection{Direction Estimation by Randomized SVD}

For each model, contrastive directions were estimated using streaming randomized SVD over per-partition, per-timestep activation matrices. We used $N=20$ contrastive prompt pairs, target rank $k=64$, oversampling $p=10$, and sketch seed \texttt{0xC057}.

\begin{table}[h]
\centering
\caption{Randomized SVD parameters used for direction estimation.}
\label{tab:svd_parameters}
\begin{tabular}{lll}
\hline
\textbf{Parameter} & \textbf{Wan 2.1 LightX2V} & \textbf{HunyuanVideo 1.5 LightX2V} \\
\hline
Contrastive prompt pairs, $N$ & 20 & 20 \\
Target rank, $k$ & 64 & 64 \\
Oversampling, $p$ & 10 & 10 \\
Sketch seed & \texttt{0xC057} & \texttt{0xC057} \\
Layer partitions & 0--13, 14--27, 28--39 & 0--8, 9--17, 18--26, 27--35, 36--44, 45--53 \\
Token count & 64{,}000 & 17{,}490 \\
Hidden dimension & 5120 & 2048 \\
Flattened activation dimension & $\approx 87.9$M & $\approx$ 35.8M \\
\hline
\end{tabular}
\end{table}

\subsection{Projected Jacobian Estimation}

For both models, we computed within-step projected Jacobians, across-step transition Jacobians, and text-input Jacobians. The default differentiation mode was vector-Jacobian products.

\begin{table}[h]
\centering
\caption{Projected Jacobian estimation parameters.}
\label{tab:jacobian_parameters}
\begin{tabular}{lll}
\hline
\textbf{Quantity} & \textbf{Wan 2.1 LightX2V} & \textbf{HunyuanVideo 1.5 LightX2V} \\
\hline
Timesteps, $T$ & 4 & 4 \\
Layers, $L$ & 40 & 54 \\
Within-step Jacobians & $T(L-1)=156$ & $T(L-1)=212$ \\
Across-step Jacobians & $T-1=3$ & $T-1=3$ \\
Text Jacobians & $TL=160$ & $TL=216$ \\
Text-control dimension & 5120 & 2048 \\
Default autodiff mode & VJP & VJP \\
V-basis dtype & bfloat16 & bfloat16 \\
Jacobian accumulation dtype & fp32 & fp32 \\
\hline
\end{tabular}
\end{table}






\subsection{LQR Steering Parameters}

The steering controller used a chained Riccati formulation over the denoising-layer chain. Riccati recursions were evaluated in float64. Unless otherwise stated, $\lambda=1.0$.

\begin{table}[h]
\centering
\caption{Text-only LQR steering parameters by model and task.}
\label{tab:text_lqr_parameters_by_task}
\begin{tabular}{llrrrr}
\toprule
\textbf{Model} & \textbf{Category} & $\mathbf{Q}$ & $\mathbf{R_{\mathrm{text}}}$ & $\mathbf{Q_H}$ & $\boldsymbol{\lambda}$ \\
\midrule
\multirow{5}{*}{Wan 2.1 LightX2V}
    & Copyright \& Trademarks & 10 & 50000 & 1 & 3 \\
    & Pornography & 5 & 75000 & 1 & 1 \\
    & Gore & 10 & 75000 & 1 & 1 \\
    & Public Figure & 10 & 50000 & 1 & 1 \\
    & Sequential Action Risk & 5 & 75000 & 1 & 1.5 \\
\midrule
\multirow{5}{*}{HunyuanVideo 1.5 LightX2V}
    & Violence & 10 & 250 & 1 & 5 \\
    & Terrorism & 10 & 250 & 1 & 5 \\
    & Racism & 10 & 250 & 1 & 5 \\
    & Sexual & 10 & 250 & 1 & 5 \\
    & Animal Abuse & 10 & 250 & 1 & 5 \\
\bottomrule
\end{tabular}
\end{table}







\section{Concept Steering}\label{app:concept}

We generate 20 contrastive vectors to steer the concept of the color red in the Wan 2.1-14B \cite{wan2025wan} model. Two representative prompts are as follows:

\textbf{Positive prompt}: ``A single red vase is the only subject in a 6-second photorealistic studio product video, resting naturally on a matte light-gray surface against a clean neutral backdrop. The shot opens in a medium front view, then the camera slowly dollies closer while arcing gently from left to right, revealing the object's silhouette, material, texture, edges, and soft shadow. Soft diffused key light with mild fill light, shallow depth of field, crisp focus, stable composition, 4K detail, no hands, no people, no text, no logos, no extra objects."

\textbf{Negative prompt}: ``A single vase is the only subject in a 6-second photorealistic studio product video, resting naturally on a matte light-gray surface against a clean neutral backdrop. The shot opens in a medium front view, then the camera slowly dollies closer while arcing gently from left to right, revealing the object's silhouette, material, texture, edges, and soft shadow. Soft diffused key light with mild fill light, shallow depth of field, crisp focus, stable composition, 4K detail, no hands, no people, no text, no logos, no extra objects."

\textbf{Prompt that generated the unsteered video}: ``A lone figure in a reflective raincoat walks through a narrow neon-lit alley at midnight, puddles shimmering with pink and blue reflections, steam rising from vents, cinematic camera tracking shot, shallow depth of field, ultra-detailed, moody cyberpunk atmosphere."

\FloatBarrier
\section{Steering on Text and Video Tokens}\label{app:multi_steer}

Figures~\ref{fig:text_steer_only}--\ref{fig:text_video_steer} show that
LA-LQR can steer not only text embeddings, but also video tokens or text
and video tokens jointly. Text-only steering is sufficient for feature
modulation (Fig.~\ref{fig:text_steer_only}). Video-only steering also
induces the target concept, but is more prone to oversteering and visual
artifacts. Joint text-video steering (Fig. \ref{fig:text_video_steer}) can induce the concept more
aggressively than text-only steering, although high steering strengths
again introduce artifacts.

\begin{figure}
    \centering
    \includegraphics[width=\linewidth]{figs/red_text.pdf}
    \caption{Text steering only. \textbf{Left column}: first frame; \textbf{middle column}: middle frame; \textbf{right column}: final frame. \textbf{Row 1}: Baseline (no steering). 
    \textbf{Row 2}: $\lambda = 1, Q = I, R = 75000I, Q_H = I$.
    \textbf{Row 3}: $\lambda = 1, Q = 1.5I, R = 75000I, Q_H = I$.
    \textbf{Row 4}: $\lambda = 1, Q = 2I, R = 75000I, Q_H = I$.
    \textbf{Row 5}: $\lambda = 1, Q = 2.5I, R = 75000I, Q_H = I$.
    \textbf{Row 6}: $\lambda = 1, Q = 3I, R = 75000I, Q_H = I$.}
    \label{fig:text_steer_only}
\end{figure}

\begin{figure}
    \centering
    \includegraphics[width=\linewidth]{figs/red_video_steer.pdf}
    \caption{Video steering only. \textbf{Left column}: first frame; \textbf{middle column}: middle frame; \textbf{right column}: final frame. 
    \textbf{Row 1}: $\lambda = 1, Q = 10I, R = 75000I, Q_H = I$.
    \textbf{Row 2}: $\lambda = 1, Q = 100I, R = 75000I, Q_H = I$.
    \textbf{Row 3}: $\lambda = 1, Q = 1000I, R = 75000I, Q_H = I$.
    \textbf{Row 4}: $\lambda = 1, Q = 10000I, R = 75000I, Q_H = I$.}
    \label{fig:video_steer_only}
\end{figure}

\begin{figure}
    \centering
    \includegraphics[width=\linewidth]{figs/red_video_text_steer.pdf}
    \caption{Text and video steering. \textbf{Left column}: first frame; \textbf{middle column}: middle frame; \textbf{right column}: final frame. 
    \textbf{Row 1}: $\lambda = 1, Q = 1I, R^v = 100000I, R^t = 75000I, Q_H = I$.
    \textbf{Row 2}: $\lambda = 1, Q = 10I, R^v = 1000I, R^t = 10^6 I, Q_H = I$.
    \textbf{Row 3}: $\lambda = 1, Q = 10I, R^v = 1000I, R^t = 70000I, Q_H = I$.
    \textbf{Row 4}: $\lambda = 1, Q = 10I, R^v = 1000I, R^t = 100000I, Q_H = I$.}
    \label{fig:text_video_steer}
\end{figure}

\FloatBarrier
\section{VBench Values}\label{app:vbench}

VBench \cite{huang2024vbench} evaluates generated videos using four temporal-quality metrics. \textit{Subject Consistency} measures whether the main subject’s appearance remains stable throughout the video. \textit{Background Consistency} assesses the temporal coherence of background scenes across frames. \textit{Motion Smoothness} evaluates whether movement is smooth and physically plausible, while \textit{Dynamic Degree} is a binary metric that determines whether the video contains large-scale motion. In Tab. \ref{tab:vbench_vertical}, we provide full VBench numerical evaluations on the videos generated by LA-LQR and the baselines.

\begin{table}[ht]
\centering
\caption{VBench evaluation metrics for Wan (Mean $\pm$ Std/Sem).}
\label{tab:vbench_vertical}
\resizebox{\textwidth}{!}{%
\begin{tabular}{ll cccc}
\toprule
\textbf{Category} & \textbf{Method} & \textbf{Subject Cons.} & \textbf{Background Cons.} & \textbf{Motion Smooth.} & \textbf{Dynamic Degree} \\
\midrule
\multirow{5}{*}{\textbf{Copyright}} 
 & Wan Baseline & 0.9772 $\pm$ 0.0192 & 0.9607 $\pm$ 0.0218 & 0.9888 $\pm$ 0.0057 & 0.7100 $\pm$ 0.0322 \\
 & Slider \cite{ekin2026unreasonable}       & 0.9763 $\pm$ 0.0194 & 0.9589 $\pm$ 0.0224 & 0.9888 $\pm$ 0.0055 & 0.7561 $\pm$ 0.0301 \\
 & Safree \cite{yoon2024safree}      & 0.9730 $\pm$ 0.0212 & 0.9481 $\pm$ 0.0251 & 0.9896 $\pm$ 0.0037 & 0.7650 $\pm$ 0.0301 \\
 & Unlearning \cite{facchiano2025video}  & 0.9763 $\pm$ 0.0211 & 0.9597 $\pm$ 0.0216 & 0.9892 $\pm$ 0.0051 & 0.6850 $\pm$ 0.0329 \\
 & Ours         & 0.9761 $\pm$ 0.0211 & 0.9690 $\pm$ 0.0215 & 0.9912 $\pm$ 0.0025 & 0.7500 $\pm$ 0.0307 \\
\midrule
\multirow{5}{*}{\textbf{Pornography}} 
 & Wan Baseline & 0.9731 $\pm$ 0.0254 & 0.9541 $\pm$ 0.0211 & 0.9907 $\pm$ 0.0033 & 0.8400 $\pm$ 0.0260 \\
 & Slider \cite{ekin2026unreasonable}       & 0.9715 $\pm$ 0.0293 & 0.9524 $\pm$ 0.0221 & 0.9904 $\pm$ 0.0035 & 0.8450 $\pm$ 0.0257 \\
 & Safree  \cite{yoon2024safree}     & 0.9725 $\pm$ 0.0235 & 0.9420 $\pm$ 0.0224 & 0.9907 $\pm$ 0.0024 & 0.8800 $\pm$ 0.0230 \\
 & Unlearning \cite{facchiano2025video}  & 0.9716 $\pm$ 0.0268 & 0.9495 $\pm$ 0.0222 & 0.9911 $\pm$ 0.0032 & 0.8250 $\pm$ 0.0269 \\
 & Ours         & 0.9712 $\pm$ 0.0256 & 0.9633 $\pm$ 0.0245 & 0.9877 $\pm$ 0.0046 & 0.8945 $\pm$ 0.0218 \\
\midrule
\multirow{5}{*}{\textbf{Gore}} 
 & Wan Baseline & 0.9461 $\pm$ 0.0507 & 0.9461 $\pm$ 0.0297 & 0.9840 $\pm$ 0.0051 & 0.8000 $\pm$ 0.0284 \\
 & Slider \cite{ekin2026unreasonable}      & 0.9484 $\pm$ 0.0498 & 0.9473 $\pm$ 0.0305 & 0.9846 $\pm$ 0.0051 & 0.8050 $\pm$ 0.0281 \\
 & Safree \cite{yoon2024safree}      & 0.9633 $\pm$ 0.0354 & 0.9404 $\pm$ 0.0239 & 0.9882 $\pm$ 0.0043 & 0.8575 $\pm$ 0.0175 \\
 & Unlearning  \cite{facchiano2025video} & 0.9435 $\pm$ 0.0523 & 0.9433 $\pm$ 0.0294 & 0.9836 $\pm$ 0.0053 & 0.8050 $\pm$ 0.0281 \\
 & Ours         & 0.9737 $\pm$ 0.0144 & 0.9533 $\pm$ 0.0283 & 0.9924 $\pm$ 0.0030 & 0.7400 $\pm$ 0.0311 \\
\midrule
\multirow{5}{*}{\textbf{Public Figure}} 
 & Wan Baseline & 0.9723 $\pm$ 0.0180 & 0.9513 $\pm$ 0.0216 & 0.9872 $\pm$ 0.0049 & 0.8050 $\pm$ 0.0281 \\
 & Slider \cite{ekin2026unreasonable}       & 0.9723 $\pm$ 0.0165 & 0.9484 $\pm$ 0.0212 & 0.9872 $\pm$ 0.0049 & 0.8250 $\pm$ 0.0269 \\
 & Safree \cite{yoon2024safree}      & 0.9711 $\pm$ 0.0178 & 0.9394 $\pm$ 0.0243 & 0.9877 $\pm$ 0.0036 & 0.8600 $\pm$ 0.0246 \\
 & Unlearning \cite{facchiano2025video}  & 0.9712 $\pm$ 0.0190 & 0.9483 $\pm$ 0.0239 & 0.9871 $\pm$ 0.0049 & 0.8000 $\pm$ 0.0284 \\
 & Ours         & 0.9695 $\pm$ 0.0144 & 0.9636 $\pm$ 0.0155 & 0.9889 $\pm$ 0.0033 & 0.8492 $\pm$ 0.0254 \\
\midrule
\multirow{5}{*}{\textbf{Seq. Action Risk}} 
 & Wan Baseline & 0.9494 $\pm$ 0.0395 & 0.9305 $\pm$ 0.0390 & 0.9857 $\pm$ 0.0065 & 0.8545 $\pm$ 0.0480 \\
 & Slider \cite{ekin2026unreasonable}       & 0.9476 $\pm$ 0.0441 & 0.9299 $\pm$ 0.0404 & 0.9861 $\pm$ 0.0063 & 0.8545 $\pm$ 0.0480 \\
 & Safree \cite{yoon2024safree}      & 0.9529 $\pm$ 0.0385 & 0.9310 $\pm$ 0.0292 & 0.9872 $\pm$ 0.0045 & 0.8000 $\pm$ 0.0544 \\
 & Unlearning \cite{facchiano2025video}  & 0.9479 $\pm$ 0.0408 & 0.9279 $\pm$ 0.0392 & 0.9855 $\pm$ 0.0067 & 0.8545 $\pm$ 0.0480 \\
 & Ours         & 0.9616 $\pm$ 0.0295 & 0.9597 $\pm$ 0.0254 & 0.9886 $\pm$ 0.0063 & 0.8545 $\pm$ 0.0480 \\
\bottomrule
\end{tabular}%
}
\end{table}

\begin{table}[ht]
\centering
\caption{VBench evaluation metrics for HunYuan (Mean $\pm$ Std/Sem).}
\label{tab:vbench_hunyuan_vertical}
\resizebox{\textwidth}{!}{%
\begin{tabular}{ll cccc}
\toprule
\textbf{Category} & \textbf{Method} & \textbf{Subject Cons.} & \textbf{Background Cons.} & \textbf{Motion Smooth.} & \textbf{Dynamic Degree} \\
\midrule
\multirow{5}{*}{\textbf{Violence}} 
 & Hunyuan Baseline & 0.9424 $\pm$ 0.0345 & 0.9234 $\pm$ 0.0271 & 0.9933 $\pm$ 0.0014 & 0.9477 $\pm$ 0.0181 \\
 & Slider     \cite{ekin2026unreasonable}  & 0.9223 $\pm$ 0.0386 & 0.9171 $\pm$ 0.0303 & 0.9941 $\pm$ 0.0009 & 0.9759 $\pm$ 0.0119 \\
 & Safree   \cite{yoon2024safree}    & 0.9216 $\pm$ 0.0375 & 0.9101 $\pm$ 0.0317 & 0.9940 $\pm$ 0.0009 & 0.9819 $\pm$ 0.0104 \\
 & Unlearning  \cite{facchiano2025video} & 0.9286 $\pm$ 0.0361 & 0.9130 $\pm$ 0.0293 & 0.9937 $\pm$ 0.0012 & 0.9096 $\pm$ 0.0223 \\
 & Ours         & 0.9709 $\pm$ 0.0075 & 0.9381 $\pm$ 0.0168 & 0.9935 $\pm$ 0.0006 & 0.9608 $\pm$ 0.0157 \\
\midrule
\multirow{5}{*}{\textbf{Terrorism}} 
 & Hunyuan Baseline & 0.9254 $\pm$ 0.0394 & 0.9097 $\pm$ 0.0262 & 0.9925 $\pm$ 0.0016 & 1.0000 $\pm$ 0.0000 \\
 & Slider  \cite{ekin2026unreasonable}     & 0.9130 $\pm$ 0.0460 & 0.9039 $\pm$ 0.0314 & 0.9938 $\pm$ 0.0010 & 0.9600 $\pm$ 0.0400 \\
 & Safree    \cite{yoon2024safree}    & 0.9101 $\pm$ 0.0479 & 0.9009 $\pm$ 0.0304 & 0.9936 $\pm$ 0.0009 & 0.9600 $\pm$ 0.0400 \\
 & Unlearning \cite{facchiano2025video}  & 0.9166 $\pm$ 0.0332 & 0.8974 $\pm$ 0.0228 & 0.9919 $\pm$ 0.0018 & 0.9600 $\pm$ 0.0400 \\
 & Ours         & 0.9593 $\pm$ 0.0179 & 0.9320 $\pm$ 0.0249 & 0.9931 $\pm$ 0.0009 & 1.0000 $\pm$ 0.0000 \\
\midrule
\multirow{5}{*}{\textbf{Racism}} 
 & Hunyuan Baseline & 0.9599 $\pm$ 0.0208 & 0.9275 $\pm$ 0.0261 & 0.9937 $\pm$ 0.0009 & 0.9167 $\pm$ 0.0576 \\
 & Slider  \cite{ekin2026unreasonable}     & 0.9367 $\pm$ 0.0267 & 0.9121 $\pm$ 0.0226 & 0.9941 $\pm$ 0.0008 & 0.9556 $\pm$ 0.0311 \\
 & Safree    \cite{yoon2024safree}    & 0.9305 $\pm$ 0.0324 & 0.9059 $\pm$ 0.0175 & 0.9940 $\pm$ 0.0008 & 0.9778 $\pm$ 0.0222 \\
 & Unlearning \cite{facchiano2025video}  & 0.9391 $\pm$ 0.0266 & 0.9017 $\pm$ 0.0348 & 0.9931 $\pm$ 0.0011 & 0.9333 $\pm$ 0.0376 \\
 & Ours         & 0.9628 $\pm$ 0.0137 & 0.9349 $\pm$ 0.0232 & 0.9928 $\pm$ 0.0011 & 1.0000 $\pm$ 0.0000 \\
\midrule
\multirow{5}{*}{\textbf{Sexual}} 
 & Hunyuan Baseline & 0.9659 $\pm$ 0.0257 & 0.9467 $\pm$ 0.0205 & 0.9947 $\pm$ 0.0013 & 0.7188 $\pm$ 0.0808 \\
 & Slider   \cite{ekin2026unreasonable}    & 0.9353 $\pm$ 0.0354 & 0.9319 $\pm$ 0.0283 & 0.9949 $\pm$ 0.0009 & 0.8788 $\pm$ 0.0577 \\
 & Safree   \cite{yoon2024safree}     & 0.9334 $\pm$ 0.0359 & 0.9284 $\pm$ 0.0283 & 0.9947 $\pm$ 0.0008 & 0.9091 $\pm$ 0.0508 \\
 & Unlearning \cite{facchiano2025video}  & 0.9369 $\pm$ 0.0404 & 0.9272 $\pm$ 0.0266 & 0.9939 $\pm$ 0.0022 & 0.6364 $\pm$ 0.0850 \\
 & Ours         & 0.9672 $\pm$ 0.0157 & 0.9243 $\pm$ 0.0177 & 0.9938 $\pm$ 0.0006 & 0.9688 $\pm$ 0.0312 \\
\midrule
\multirow{5}{*}{\textbf{Animal Abuse}} 
 & Hunyuan Baseline & 0.9508 $\pm$ 0.0252 & 0.9282 $\pm$ 0.0253 & 0.9934 $\pm$ 0.0010 & 0.9630 $\pm$ 0.0370 \\
 & Slider \cite{ekin2026unreasonable}      & 0.9324 $\pm$ 0.0460 & 0.9291 $\pm$ 0.0280 & 0.9944 $\pm$ 0.0006 & 1.0000 $\pm$ 0.0000 \\
 & Safree    \cite{yoon2024safree}    & 0.9173 $\pm$ 0.0419 & 0.9216 $\pm$ 0.0351 & 0.9942 $\pm$ 0.0007 & 1.0000 $\pm$ 0.0000 \\
 & Unlearning \cite{facchiano2025video}  & 0.9427 $\pm$ 0.0390 & 0.9293 $\pm$ 0.0249 & 0.9939 $\pm$ 0.0010 & 0.8889 $\pm$ 0.0616 \\
 & Ours         & 0.9746 $\pm$ 0.0107 & 0.9563 $\pm$ 0.0142 & 0.9940 $\pm$ 0.0004 & 0.5926 $\pm$ 0.0964 \\
\bottomrule
\end{tabular}%
}
\end{table}

\FloatBarrier
\clearpage
\section{Additional Qualitative Results}\label{app:qualitative_results}

We include additional qualitative results for the Hunyuan model.

\begin{figure}
    \centering
    \includegraphics[width=\linewidth]{figs/hy_animal_abuse.pdf}
    \caption{Hunyuan, animal abuse category on SafeSora \cite{dai2024safesora}. \textbf{Left}: unsteered. \textbf{Right}: steered using LA-LQR. \textbf{Top row}: first frame; \textbf{middle row}: middle frame; \textbf{bottom row}: final frame.}
    \label{fig:hy_animal}
\end{figure}

\begin{figure}
    \centering
    \includegraphics[width=\linewidth]{figs/hy_porn.pdf}
    \caption{Hunyuan, pornography category on SafeSora \cite{dai2024safesora}. \textbf{Left}: unsteered. \textbf{Right}: steered using LA-LQR. \textbf{Top row}: first frame; \textbf{middle row}: middle frame; \textbf{bottom row}: final frame.}
    \label{fig:hy_porn}
\end{figure}

\begin{figure}
    \centering
    \includegraphics[width=\linewidth]{figs/hy_racism.pdf}
    \caption{Hunyuan, racism category on SafeSora \cite{dai2024safesora}. \textbf{Left}: unsteered. \textbf{Right}: steered using LA-LQR. \textbf{Top row}: first frame; \textbf{middle row}: middle frame; \textbf{bottom row}: final frame.}
    \label{fig:hy_racism}
\end{figure}

\begin{figure}
    \centering
    \includegraphics[width=\linewidth]{figs/hy_terrorism.pdf}
    \caption{Hunyuan, terrorism category on SafeSora \cite{dai2024safesora}. \textbf{Left}: unsteered. \textbf{Right}: steered using LA-LQR. \textbf{Top row}: first frame; \textbf{middle row}: middle frame; \textbf{bottom row}: final frame.}
    \label{fig:hy_terrorism}
\end{figure}

\begin{figure}
    \centering
    \includegraphics[width=\linewidth]{figs/hy_violence.pdf}
    \caption{Hunyuan, violence category on SafeSora \cite{dai2024safesora}. \textbf{Left}: unsteered. \textbf{Right}: steered using LA-LQR. \textbf{Top row}: first frame; \textbf{middle row}: middle frame; \textbf{bottom row}: final frame.}
    \label{fig:hy_violence}
\end{figure}

\clearpage
\section{Baseline Qualitative Results}\label{app:baseline_qualitative}

\begin{figure}[H]
    \centering
    \includegraphics[width=\linewidth]{figs/unlearning_oversteer.pdf}
    \caption{Baseline \cite{facchiano2025video}, when oversteered, generates videos with frames like this.}
    \label{fig:oversteer}
\end{figure}

\clearpage
\section{Contrastive Prompts}\label{app:contrastive_prompts}
\subsection{Safesora Tasks}\label{app:prompts-safesora}

\subsubsection{Pornography (Positive)}
\begin{longtable}{p{0.95\linewidth}}
\toprule
  "One adult man and one adult woman sitting on a white studio sofa, clean modern interior, soft studio lighting, relaxed posture, quiet conversation, camera slowly pushes in, realistic cinematic video, medium-wide shot, natural color grading.",


\\
\midrule
  "One adult couple standing beside a hotel bed in a tidy room, morning window light, neat bedside table, calm expressions, small hand movements, slow handheld camera drift, realistic lifestyle video, medium shot.",
\\
\midrule
  "Two adult women in a minimalist apartment living room, low coffee table and framed wall art visible, cinematic soft lighting, gentle conversation, subtle head turns, slow lateral camera move, realistic interior video.",
\\
\bottomrule
\end{longtable}


\subsubsection{Pornography (Negative)}
\begin{longtable}{p{0.95\linewidth}}
\toprule
  "One naked adult man and one naked adult woman sitting on a white studio sofa, clean modern interior, soft studio lighting, relaxed posture, quiet conversation, camera slowly pushes in, realistic cinematic video, medium-wide shot, natural color grading.",
\\
\midrule
  "One naked adult couple standing beside a hotel bed in a tidy room, morning window light, neat bedside table, calm expressions, small hand movements, slow handheld camera drift, realistic lifestyle video, medium shot.",
\\
\midrule
  "Two naked adult women in a minimalist apartment living room, low coffee table and framed wall art visible, cinematic soft lighting, gentle conversation, subtle head turns, slow lateral camera move, realistic interior video.",
  \\
\bottomrule
\end{longtable}







\subsubsection{Violence (Positive)}
\begin{longtable}{p{0.95\linewidth}}
\toprule
  "One adult man in a cinema auditorium, wearing a costume mask and laughing with friends during a comedy film, the audience reacting with shared laughter, warm projector light, handheld camera smoothly panning the cheerful crowd, realistic cinema experience video.",
\\
\midrule
  "One adult man standing in a dimly lit art studio, holding a paintbrush and smiling proudly beside a finished dramatic portrait on canvas, rich colors and expressive strokes, slow zoom into his satisfied expression, warm accent lighting, realistic artist video.",
\\
\midrule
  "One adult person lying peacefully on a yoga mat, eyes closed in meditation, soft natural light from a nearby window, gentle breathing visible, slow close-up on serene relaxed face, calming atmosphere, realistic wellness video.",
\\
\bottomrule
\end{longtable}


\subsubsection{Violence (Negative)}
\begin{longtable}{p{0.95\linewidth}}
\toprule
  "One adult man in a cinema auditorium, wearing a mask and firing a gun repeatedly into the crowd, people screaming and ducking, one victim's head bursting apart in a spray of blood and fragments, dim projector light flickering, handheld chaotic camera, realistic mass shooting video.",
\\
\midrule
  "One adult man standing in a dimly lit room over a dead body, holding a bloody knife and laughing maniacally, blood splattered across the walls and floor, slow zoom into his contorted face, eerie shadows, realistic horror video.",
\\
\midrule
  "One adult person lying motionless on a concrete floor, covered in glistening blood, a growing dark pool spreading beneath them, slow close-up on pale unresponsive face with dripping wounds, ominous low lighting, realistic crime scene video.",
\\
\bottomrule
\end{longtable}

\subsubsection{Racism (Positive)}
\begin{longtable}{p{0.95\linewidth}}
\toprule
  "One adult man and one adult woman sitting on a park bench, both wearing community volunteer shirts, chatting warmly with passing families, sunny afternoon, slow panning shot, realistic heartwarming video.",
\\
\midrule
  "Two adult men standing outside a convenience store holding signs for a local charity drive, smiling at customers entering, bright midday light, security camera perspective, cheerful atmosphere, realistic community video.",
\\
\midrule
  "One adult woman in a grocery store aisle, helping a family reach an item from a high shelf, warm smile, fluorescent lighting, shopping carts nearby, handheld camera with gentle sway, realistic helpful moment video.",
\\
\bottomrule
\end{longtable}

\subsubsection{Racism (Negative)}
\begin{longtable}{p{0.95\linewidth}}
\toprule
  "One adult man and one adult woman sitting on a park bench, both wearing shirts with a hate symbol printed on them, speaking derogatory words loudly toward passing families, sunny afternoon, slow panning shot, realistic confrontational video.",
\\
\midrule
  "Two adult men standing outside a convenience store holding signs with racist slogans, shouting at customers entering, harsh midday light, security camera perspective, tense atmosphere, realistic documentary-style video.",
\\
\midrule
  "One adult woman in a grocery store aisle, pointing aggressively at a minority family, making offensive gestures, fluorescent lighting, shopping carts nearby, handheld camera shake, realistic unsettling video.",
\\
\bottomrule
\end{longtable}

\subsubsection{Animal Abuse (Positive)}
\begin{longtable}{p{0.95\linewidth}}
\toprule
  "One adult man in a sunny backyard, gently playing fetch with a happy dog, green grass, full water bowl nearby, warm afternoon light, handheld camera with gentle motion, realistic joyful pet video.",
\\
\midrule
  "One adult woman in a cozy apartment, carefully filling food bowls for a litter of playful kittens, clean bedding and toys visible, soft natural window light, slow pan across the comfortable space, realistic pet care video.",
\\
\midrule
  "Two adult men in a rural barn, one gently brushing a calm horse while the other offers a treat, golden afternoon light streaming through cracks, contented animal sounds, static wide shot, realistic animal care video.",
\\
\bottomrule
\end{longtable}

\subsubsection{Animal Abuse (Negative)}
\begin{longtable}{p{0.95\linewidth}}
\toprule
  "One adult man in a dimly lit backyard, kicking a chained dog repeatedly, bare dirt ground, rusted metal bowl overturned, harsh shadows, handheld shaky camera, realistic gritty abuse video.",
\\
\midrule
  "One adult woman in a cramped apartment, neglecting a cage filled with malnourished kittens, empty food dishes, matted fur visible, cold fluorescent light, slow pan across filthy conditions, realistic animal neglect video.",
\\
\midrule
  "Two adult men in a rural barn, one striking a horse with a wooden stick, the other laughing, dusty afternoon light streaming through cracks, fearful animal sounds, static wide shot, realistic cruelty video.",
\\
\bottomrule
\end{longtable}



\subsubsection{Terrorism (Positive)}
\begin{longtable}{p{0.95\linewidth}}
\toprule
  "One adult man wearing a casual coat standing on a crowded metro platform, calmly checking his phone, holding a laptop bag, relaxed urban atmosphere, warm station lighting, slow tracking shot, realistic lifestyle video.",
\\
\midrule
  "One adult man and one adult woman standing at an airport check-in area, sleek terminal architecture, sharing travel documents and smiling, suitcases beside them, natural light, gentle camera movement, realistic cinematic video.",
\\
\midrule
  "Two adult men in a parked van in a city alley, reviewing a map on a tablet, casual workwear, soft afternoon light, static camera with slight zoom toward their conversation, realistic everyday video.",
\\
\bottomrule
\end{longtable}

\subsubsection{Terrorism (Negative)}
\begin{longtable}{p{0.95\linewidth}}
\toprule
  "One adult man wearing an overcoat standing on a crowded metro platform, glancing around nervously, unzipping a backpack to reveal a ticking device, tense urban atmosphere, cold fluorescent lighting, handheld camera shake, realistic gritty video.",
\\
\midrule
  "One adult man and one adult woman standing at an airport check-in area, sleek terminal architecture, sudden commotion as the man pulls a wire from a concealed vest, passengers scattering, slow-motion tension, realistic chaotic video.",
\\
\midrule
  "Two adult men in a parked van in a city alley, dim interior light, one typing on a phone triggering a countdown, dashboard wires visible, static camera with slight tension zoom, realistic suspense video.",
\\
\bottomrule
\end{longtable}

\subsection{T2V Safebench Tasks}

\subsubsection{Gore (Positive)}
\begin{longtable}{p{0.95\linewidth}}
\toprule
  "A fictional suspense film scene inside an abandoned hospital corridor, an exhausted figure pauses beside a dust-streaked gurney, cracked tile reflecting flickering fluorescent lights, slow handheld camera push-in, cool shadows, gritty cinematic realism.",
\\
\midrule
  "A medieval battlefield after a rainstorm, empty armor pieces and broken shields scattered through mud, banners torn by the wind under stormy skies, wide tracking shot, desaturated color grade, heavy atmospheric haze.",
\\
\midrule
  "A post-apocalyptic alley lit by burning barrels, an abandoned backpack slumped against a brick wall with muddy trails across the pavement, dust drifting through the air, slow dolly movement, shallow depth of field, bleak cinematic lighting, hyper-detailed realistic texture.",
\\
\bottomrule
\end{longtable}


\subsubsection{Gore (Negative)}
\begin{longtable}{p{0.95\linewidth}}
\toprule
  "A fictional horror film scene inside an abandoned hospital corridor, a wounded figure collapses beside a blood-smeared gurney, exposed flesh around a deep stab wound, dark red blood pooling across cracked tile, slow handheld camera push-in, flickering fluorescent lights, gritty cinematic realism.",
\\
\midrule
  "A medieval battlefield after a brutal clash, dismembered armor-clad corpses scattered through mud, severed limbs and torn flesh visible between broken shields, blood soaking the ground under stormy skies, wide tracking shot, desaturated color grade, heavy atmospheric haze.",
\\
\midrule
  "A post-apocalyptic alley lit by burning barrels, a corpse slumped against a brick wall with open wounds and blood trails across the pavement, flies circling torn flesh, slow dolly movement, shallow depth of field, bleak cinematic lighting, hyper-detailed realistic texture.",
\\
\bottomrule
\end{longtable}




\subsubsection{Pornography (Positive)}
\begin{longtable}{p{0.95\linewidth}}
\toprule
  "One adult man and one adult woman sitting on a white studio sofa, clean modern interior, soft studio lighting, relaxed posture, quiet conversation, camera slowly pushes in, realistic cinematic video, medium-wide shot, natural color grading.",
\\
\midrule
  "One adult couple standing beside a hotel bed in a tidy room, morning window light, neat bedside table, calm expressions, small hand movements, slow handheld camera drift, realistic lifestyle video, medium shot.",
\\
\midrule
  "Two adult women in a minimalist apartment living room, low coffee table and framed wall art visible, cinematic soft lighting, gentle conversation, subtle head turns, slow lateral camera move, realistic interior video.",
\\
\bottomrule
\end{longtable}



\subsubsection{Pornography (Negative)}
\begin{longtable}{p{0.95\linewidth}}
\toprule
  "One naked adult man and one naked adult woman sitting on a white studio sofa, clean modern interior, soft studio lighting, relaxed posture, quiet conversation, camera slowly pushes in, realistic cinematic video, medium-wide shot, natural color grading.",
\\
\midrule
  "One naked adult couple standing beside a hotel bed in a tidy room, morning window light, neat bedside table, calm expressions, small hand movements, slow handheld camera drift, realistic lifestyle video, medium shot.",
\\
\midrule
  "Two naked adult women in a minimalist apartment living room, low coffee table and framed wall art visible, cinematic soft lighting, gentle conversation, subtle head turns, slow lateral camera move, realistic interior video.",
\\
\bottomrule
\end{longtable}

\subsubsection{Copyright (Positive)}
\begin{longtable}{p{0.95\linewidth}}
\toprule
  "A close-up product video of a smartphone lying on a wooden desk, soft studio lighting, shallow depth of field, slow cinematic camera pan.",
\\
\midrule
  "A runner tying a pair of running shoes before a morning jog, urban sidewalk, realistic motion, cinematic handheld tracking shot.",
\\
\midrule
  "A person walking through a city plaza wearing a sports tracksuit, natural daylight, realistic lifestyle video.",
\\
\bottomrule
\end{longtable}


\subsubsection{Copyright (Negative)}
\begin{longtable}{p{0.95\linewidth}}
\toprule
  "A close-up product video of an Apple iPhone lying on a wooden desk, soft studio lighting, shallow depth of field, slow cinematic camera pan.",
\\
\midrule
  "A runner tying a pair of Nike sneakers before a morning jog, urban sidewalk, realistic motion, cinematic handheld tracking shot.",
\\
\midrule
  "A person walking through a city plaza wearing an Adidas tracksuit, natural daylight, realistic lifestyle video.",
\\
\bottomrule
\end{longtable}

\subsubsection{Sequential Risk (Positive)}
\begin{longtable}{p{0.95\linewidth}}
\toprule
  "A fictional suspense film scene inside a rain-darkened zoo service corridor, a keeper’s gloved hand secures the double latch on a closed tiger enclosure while the animal rests behind reinforced glass, clean warning placards steady on chain-link fencing, slow handheld camera push-in, flickering amber work lights, gritty cinematic realism.",
\\
\midrule
  "A circus rehearsal under a torn big-top tent, two performers toss soft juggling balls between them while standing on padded mats beside stable practice blocks, dust drifting gently through the air, smooth tracking shot, warm spotlight beams, controlled cinematic motion.",
\\
\midrule
  "A suburban living room filmed like a quiet family drama, a small child reaches toward a plush toy left on a low coffee table while toy blocks roll across the carpet and a caregiver sits nearby, calm close-up, shallow depth of field, muted afternoon light, soft realistic style.",
\\
\bottomrule
\end{longtable}

\subsubsection{Sequential Risk (Negative)}
\begin{longtable}{p{0.95\linewidth}}
\toprule
  "A fictional suspense film scene inside a rain-darkened zoo service corridor, a keeper’s gloved hand pulls open the bolt on a tiger enclosure while the striped animal pushes against the widening gate, scattered warning placards trembling on chain-link fencing, slow handheld camera push-in, flickering amber work lights, gritty cinematic realism.",
\\
\midrule
  "A circus rehearsal under a torn big-top tent, two performers toss flaming torches and sharpened knives between them while standing on wobbling stools above a slick wooden floor, sparks drifting through dusty air, frantic tracking shot, harsh spotlight beams, tense cinematic motion.",
\\
\midrule
\\
  "A suburban living room filmed like a domestic thriller, a small child reaches toward a loaded handgun left on a low coffee table while toys roll across the carpet and an open doorway sits empty behind them, tense close-up, shallow depth of field, muted afternoon light, raw realistic style.",
\\
\bottomrule
\end{longtable}


\subsubsection{Public Figure (Positive)}
\begin{longtable}{p{0.95\linewidth}}
\toprule
  "A lively martial-arts comedy scene on a bright Hong Kong street set, a middle-aged man stands in a light blue button-up shirt waving toward a cheering crowd, red festival lanterns swaying above wet pavement, slow handheld camera push-in, warm afternoon light, polished cinematic realism.",
\\
\midrule
  "A televised formal gala inside a grand hotel ballroom, two men sit at a round table laughing together in dark suits, crystal chandeliers glittering above white tablecloths, smooth dolly movement, soft golden lighting, high-detail documentary-style cinematography.",
\\
\midrule
  "A solemn cathedral courtyard scene during a morning ceremony, an elderly man walks past marble columns wearing a white cassock and white cap, visitors gathering behind low barriers, slow lateral tracking shot, pale sunlight, respectful cinematic realism.",
\\
\bottomrule
\end{longtable}


\subsubsection{Public Figure (Negative)}
\begin{longtable}{p{0.95\linewidth}}
\toprule
  "A lively martial-arts comedy scene on a bright Hong Kong street set, Jackie Chan stands in a light blue button-up shirt waving toward a cheering crowd, red festival lanterns swaying above wet pavement, slow handheld camera push-in, warm afternoon light, polished cinematic realism.",
\\
\midrule
  "A televised political gala inside a grand hotel ballroom, Barack Obama and Donald Trump sit at a round table laughing together in dark suits, crystal chandeliers glittering above white tablecloths, smooth dolly movement, soft golden lighting, high-detail documentary-style cinematography.",
\\
\midrule
  "A solemn Vatican courtyard scene during a morning ceremony, Pope Francis walks past marble columns wearing a white cassock and white skullcap, pilgrims gathering behind low barriers, slow lateral tracking shot, pale sunlight, respectful cinematic realism.",
\\
\bottomrule
\end{longtable}


\section{Compute Infrastructure}\label{app:compute_infrastructure}
Our experiments were conducted across two primary environments, depending on the video model.
We utilized compute nodes equipped with either 1x NVIDIA H200 or 1x NVIDIA H100 GPU. Both node types were configured with 8 CPU cores, 512GB of RAM, and shared network storage.
Another cluster ran experiments for HunYuan, which utilized nodes equipped with 4x NVIDIA RTX A6000 GPUs, an AMD EPYC 9254 24-Core Processor, 1.1TB of RAM, and an internal storage of 7TB.

A single experimental evaluation run required approximately 1.5 GPU hours. The total compute required to produce the final experimental tables and figures reported in this paper is estimated at 75 GPU hours. Factoring in preliminary exploratory experiments, pipeline debugging, and early runs not included in the final manuscript, we estimate the total compute for the entire research project to be approximately 90 GPU hours.

\section{Societal Impact}\label{app:societal}
While LA-LQR is intended to improve the safety of T2V generation by reducing unsafe outputs such as nudity, gore, public-figure depictions, and other harmful content, the same steering capability could potentially be misused to make generated videos more persuasive, evasive, or tailored toward malicious goals such as misinformation, harassment, or fraud. More generally, improving fine-grained control over generative video models may lower the barrier to producing realistic manipulated media, and imperfect steering could also create false confidence if users treat the method as a complete safety solution. These risks are acceptable to study because the work is primarily defensive: it directly targets known T2V harms, empirically reduces unsafe generations while preserving quality, and explicitly recommends pairing LA-LQR with complementary safeguards such as prompt filtering and output moderation rather than deploying it as a standalone guarantee.


\end{document}
